\chapter{Barisan Fungsi} \label{fs:chapter}

%%%%%%%%%%%%%%%%%%%%%%%%%%%%%%%%%%%%%%%%%%%%%%%%%%%%%%%%%%%%%%%%%%%%%%%%%%%%%%

\section{Konvergensi titik demi titik dan seragam}
\label{sec:puconv}

%mbxINTROSUBSECTION

\sectionnotes{1--1.5 kuliah}

Hingga saat ini, ketika membahas limit barisan, kita membahas
barisan bilangan.
Konsep yang sangat berguna dalam analisis adalah barisan fungsi.
Sebagai contoh, solusi suatu persamaan diferensial
mungkin ditemukan dengan terlebih dahulu mencari solusi-solusi hampiran.
Solusi sebenarnya kemudian merupakan semacam limit dari solusi-solusi hampiran tersebut.

Ketika membahas barisan fungsi,
bagian yang rumit adalah adanya beberapa pengertian limit.
Mari kita uraikan dua pengertian umum
mengenai limit barisan fungsi.

\subsection{Konvergensi titik demi titik}

\begin{defn}
\index{konvergensi titik demi titik}
Untuk setiap $n \in \N$,
misalkan $f_n \colon S \to \R$ merupakan fungsi.  Barisan
$\{ f_n \}_{n=1}^\infty$
\emph{\myindex{konvergen titik demi titik}}\footnote{Kecuali dinyatakan lain,
\emph{konvergen} pada umumnya berarti \emph{konvergen titik demi titik}.}
ke $f \colon S \to \R$ jika untuk setiap $x
\in S$ berlaku
\begin{equation*}
f(x) =
\lim_{n\to\infty} f_n(x) .
\end{equation*}
\end{defn}

Limit barisan bilangan bersifat tunggal, sehingga jika barisan
$\{ f_n \}_{n=1}^\infty$
konvergen titik demi titik, fungsi limit $f$ bersifat tunggal.
Lazim dikatakan bahwa $f_n \colon S \to \R$
\emph{konvergen titik demi titik ke $f$ pada $T \subset S$}
untuk suatu $f \colon T \to \R$.  Dalam hal ini, yang dimaksud ialah
$f(x) = \lim_{n\to\infty} f_n(x)$ untuk setiap $x \in T$.  Dengan kata lain,
pembatasan $f_n$ pada $T$ konvergen titik demi titik ke $f$.


\begin{example}
Pada $[-1,1]$, barisan fungsi yang didefinisikan oleh $f_n(x) \coloneqq x^{2n}$
konvergen titik demi titik ke $f \colon [-1,1] \to \R$, dengan
\begin{equation*}
f(x) =
\begin{cases}
1 & \text{jika } x=-1 \text{ atau } x=1, \\
0 & \text{selain itu.}
\end{cases}
\avoidbreak
\end{equation*}
Lihat \figureref{x2nfig}.

\begin{myfigureht}
\myincludegraphics{x2nfig}{%
Grafik fungsi x kuadrat, x pangkat empat, x pangkat enam, dan
x pangkat enam belas.  Grafik-grafik tersebut berada pada daerah dari minus 1 hingga 1,
semuanya berawal di titik yang sama di sebelah kiri dan berakhir
di titik yang sama di sebelah kanan (karena semua fungsi bernilai 1
di titik-titik ujung).  Di bagian dalam interval, pangkat yang lebih tinggi membuat grafik
lebih datar dan lebih dekat ke sumbu x di tengah, tetapi kemudian grafik harus naik lebih
curam kembali ke 1 ketika mendekati ujung-ujung interval.}
\caption{Grafik $f_1$, $f_2$, $f_3$, dan $f_8$ untuk $f_n(x) \coloneqq
x^{2n}$.\label{x2nfig}}
\end{myfigureht}

Untuk melihat bahwa hal ini benar, pertama ambil $x \in (-1,1)$.  Maka
$0 \leq x^2 < 1$.
Kita telah melihat sebelumnya bahwa
\begin{equation*}
\sabs{x^{2n} - 0} = {(x^2)}^n \to 0 \quad \text{ketika} \quad n \to \infty .
\end{equation*}
Oleh karena itu, $\lim_{n\to\infty}f_n(x) = 0$.
Ketika $x = 1$ atau $x=-1$, berlaku $x^{2n} = 1$ untuk semua $n$, sehingga
$\lim_{n\to\infty}f_n(x) = 1$.

Perhatikan bahwa untuk $x$ yang tidak berada dalam $[-1,1]$, barisan
$\{ x^{2n} \}_{n=1}^\infty$ tidak konvergen.
Jadi, jika kita mendefinisikan $f_n$ pada himpunan yang lebih besar daripada $[-1,1]$,
maka $\{ f_n \}_{n=1}^\infty$
hanya akan konvergen titik demi titik pada $[-1,1]$.
\end{example}

Sering kali, fungsi diberikan sebagai deret.  Dalam hal ini, kita menggunakan
pengertian konvergensi titik demi titik untuk menentukan nilai-nilai fungsi tersebut.

\begin{example} \label{example:geomsumptconv}
Kita menulis
\begin{equation*}
\sum_{k=0}^\infty x^k
\end{equation*}
untuk menyatakan limit fungsi-fungsi
\begin{equation*}
f_n(x) \coloneqq \sum_{k=0}^n x^k .
\end{equation*}
Ketika mempelajari deret,
kita melihat bahwa pada $(-1,1)$, fungsi-fungsi $f_n$ konvergen titik demi titik ke
\begin{equation*}
\frac{1}{1-x} .
\end{equation*}

Hal yang perlu diperhatikan adalah bahwa walaupun
$\frac{1}{1-x}$ terdefinisi untuk semua $x \neq 1$ dan $f_n$
terdefinisi untuk semua $x$ (bahkan di $x=1$), konvergensi hanya terjadi pada $(-1,1)$.
Oleh karena itu, ketika kita menulis
\begin{equation*}
f(x) \coloneqq \sum_{k=0}^\infty x^k ,
\end{equation*}
yang dimaksud adalah bahwa $f$ terdefinisi pada $(-1,1)$ dan merupakan limit titik demi titik
dari jumlah-jumlah parsial tersebut.
\end{example}

\begin{example}
Misalkan $f_n(x) \coloneqq \sin(nx)$.  Maka $f_n$ tidak konvergen titik demi titik
ke fungsi mana pun pada interval mana pun.  Barisan tersebut mungkin konvergen di titik-titik tertentu, misalnya
ketika $x=0$ atau $x=\pi$.  Sebagai latihan, tunjukkan bahwa dalam setiap interval
$[a,b]$, terdapat $x$ sedemikian sehingga $\sin(xn)$ tidak memiliki limit
ketika $n$ menuju tak hingga.  Lihat \figureref{fig:nonconvsinxn}.
\begin{myfigureht}
\myincludegraphics{nonconvsinxn}{%
Grafik sejumlah gelombang sinus, dengan frekuensi yang lebih tinggi digambar dalam warna abu-abu
yang lebih muda.  Gambar tersebut tampak seolah-olah kita mencoba mengisi
suatu pita horizontal dengan menggambar garis-garis bergelombang.
Semua kurva berawal di titik yang sama pada
sumbu x di sebelah kiri sambil bergerak naik,
dan semuanya melalui titik yang sama pada sumbu x di
sebelah kanan sambil bergerak naik.
Semua kurva melalui pusat pada sumbu x, tetapi sebagian bergerak naik
dan sebagian bergerak turun.
Pola serupa berulang di titik-titik lain, tetapi hanya sebagian
garis berfrekuensi lebih tinggi yang melalui titik-titik tersebut.}
\caption{Grafik $\sin(nx)$ untuk
$n=1,2,\ldots,10$, dengan $n$ yang lebih besar digambar dalam warna abu-abu yang lebih muda.%
\label{fig:nonconvsinxn}}
\end{myfigureht}
\end{example}

Sebelum beralih ke konvergensi seragam, mari kita nyatakan kembali konvergensi
titik demi titik dengan cara lain.
Pembuktiannya diserahkan kepada pembaca---ini merupakan penerapan sederhana dari
definisi konvergensi barisan bilangan real.

\begin{prop} \label{ptwsconv:prop}
Misalkan $f_n \colon S \to \R$ dan $f \colon S \to \R$ merupakan fungsi.
Maka $\{ f_n \}_{n=1}^\infty$ konvergen titik demi titik ke $f$ jika dan hanya jika
untuk setiap $x \in S$ dan setiap $\epsilon > 0$, terdapat
$N \in \N$ sedemikian sehingga
\begin{equation*}
\babs{f_n(x)-f(x)} < \epsilon \quad \text{untuk semua } n \geq N .
\end{equation*}
\end{prop}

Hal yang penting adalah bahwa $N$ dapat bergantung pada $x$, bukan hanya pada
$\epsilon$.  Untuk setiap $x$, kita dapat memilih $N$ yang berbeda.
Jika kita dapat memilih satu $N$ untuk semua $x$, kita akan memperoleh apa yang disebut
konvergensi seragam.

\subsection{Konvergensi seragam}

\begin{defn}
\index{konvergensi seragam}
Misalkan $f_n \colon S \to \R$
dan $f \colon S \to \R$
merupakan fungsi.  Barisan $\{ f_n \}_{n=1}^\infty$
\emph{\myindex{konvergen seragam}} ke $f$ jika untuk
setiap $\epsilon > 0$, terdapat $N \in \N$ sedemikian sehingga
untuk semua $n \geq N$,
\begin{equation*}
\babs{f_n(x) - f(x)} < \epsilon \qquad \text{untuk semua } x \in S.
\end{equation*}
\end{defn}
\begin{myfigureht}
\myincludegraphics{uniformconv}{%
Grafik suatu fungsi yang diberi label f, dengan pita di sekeliling fungsi
yang batas-batasnya diberi label f dikurangi epsilon dan f ditambah epsilon.  Fungsi
yang diberi label f subskrip n ditampilkan sebagai garis bergelombang yang seluruhnya berada
di dalam pita tersebut.}
\caption{Dalam konvergensi seragam,
untuk $n \geq N$,
fungsi-fungsi $f_n$ berada di dalam pita $\pm\epsilon$ di sekitar $f$.%
\label{fig:uniformconv}}
\end{myfigureht}

Dalam konvergensi seragam, $N$ tidak boleh bergantung pada $x$.  Diberikan $\epsilon > 0$,
kita harus menemukan $N$ yang berlaku untuk semua $x \in S$.  Lihat
\figureref{fig:uniformconv} untuk ilustrasinya.
Konvergensi seragam
menyiratkan konvergensi titik demi titik, dan pembuktiannya mengikuti
\propref{ptwsconv:prop}:

\begin{prop}
Misalkan $\{ f_n \}_{n=1}^\infty$ merupakan barisan fungsi $f_n \colon S \to \R$.
Jika $\{ f_n \}_{n=1}^\infty$ konvergen
seragam ke $f \colon S \to \R$, maka $\{ f_n \}_{n=1}^\infty$ konvergen titik demi titik ke $f$.
\end{prop}

Pernyataan sebaliknya tidak berlaku.

\begin{example}
Fungsi-fungsi $f_n(x) \coloneqq x^{2n}$ tidak konvergen seragam pada $[-1,1]$,
meskipun konvergen titik demi titik.  Untuk melihatnya, andaikan demi kontradiksi
bahwa konvergensinya seragam.  Untuk $\epsilon \coloneqq \nicefrac{1}{2}$, harus
terdapat $N$ sedemikian sehingga $x^{2N} = \sabs{x^{2N} - 0} < \nicefrac{1}{2}$ untuk semua $x \in
(-1,1)$ (karena $f_n(x)$ konvergen ke 0 pada $(-1,1)$).  Namun, itu berarti bahwa
untuk setiap barisan $\{ x_k \}_{k=1}^\infty$ dalam $(-1,1)$ sedemikian sehingga $\lim_{k\to\infty} x_k = 1$,
berlaku $x_k^{2N} < \nicefrac{1}{2}$ untuk semua $k$.  Di sisi lain,
$x^{2N}$ merupakan fungsi kontinu terhadap $x$ (fungsi tersebut adalah polinom).  Oleh karena itu,
kita memperoleh kontradiksi
\begin{equation*}
1 = 1^{2N}  = \lim_{k\to\infty} x_k^{2N} \leq \nicefrac{1}{2} .
\end{equation*}

Namun, jika kita membatasi domain pada $[-a,a]$ dengan $0 < a < 1$, maka
$\{ f_n \}_{n=1}^\infty$ konvergen seragam ke 0 pada $[-a,a]$.  Perhatikan bahwa
$a^{2n} \to 0$ ketika $n \to \infty$.  Diberikan $\epsilon > 0$,
pilih $N \in \N$ sedemikian sehingga
$a^{2n} < \epsilon$ untuk semua $n \geq N$.
Jika $x \in [-a,a]$, maka
$\sabs{x} \leq a$.  Jadi,
untuk semua $n \geq N$ dan semua $x \in [-a,a]$,
\begin{equation*}
\sabs{x^{2n}} = \sabs{x}^{2n} \leq a^{2n} < \epsilon .
\end{equation*}
\end{example}

\subsection{Konvergensi dalam norma seragam}

Untuk fungsi terbatas, ada cara lain yang lebih abstrak untuk
memahami konvergensi seragam.  Kepada setiap fungsi terbatas, kita kaitkan
suatu bilangan tak negatif yang
mengukur \myquote{jarak} fungsi tersebut dari fungsi konstan
$0$.
Bilangan ini memungkinkan kita untuk
\myquote{mengukur}
seberapa jauh dua fungsi satu sama lain.  Kemudian kita menerjemahkan
pernyataan tentang konvergensi seragam menjadi pernyataan tentang suatu
barisan bilangan real yang konvergen ke nol.

\begin{defn} \label{def:unifnorm}
Misalkan $f \colon S \to \R$ merupakan fungsi terbatas.  Definisikan
\glsadd{not:uniformnorm}
\begin{equation*}
\snorm{f}_S \coloneqq
\sup \Bigl\{ \babs{f(x)} : x \in S \Bigr\} .
\avoidbreak
\end{equation*}
Kita menyebut $\snorm{\cdot}_S$ sebagai \emph{\myindex{norma seragam}}.
Terkadang digunakan notasi lain\footnote{Notasi dan terminologinya
belum sepenuhnya baku.  Norma ini juga disebut
\emph{\myindex{norma supremum}} atau
\emph{\myindex{norma tak hingga}}; selain
$\snorm{f}_u$ dan $\snorm{f}_S$, norma ini terkadang ditulis
sebagai $\snorm{f}_{\infty}$ atau $\snorm{f}_{\infty,S}$.},
misalnya $\snorm{f}_u$.
\end{defn}

Subskripnya adalah himpunan tempat supremum diambil.
Jadi, jika $K \subset S$, maka
\begin{equation*}
\snorm{f}_K =
\sup \Bigl\{ \babs{f(x)} : x \in K \Bigr\} .
\end{equation*}

\begin{prop}
Barisan fungsi terbatas $f_n \colon S \to \R$ konvergen
seragam ke $f \colon S \to \R$ jika dan hanya jika
\begin{equation*}
\lim_{n\to\infty} \snorm{f_n - f}_S = 0 .
\end{equation*}
\end{prop}

\begin{proof}
Pertama, andaikan
$\lim_{n\to\infty} \snorm{f_n - f}_S = 0$.  Misalkan diberikan
$\epsilon > 0$.  Terdapat $N$ sedemikian sehingga
untuk $n \geq N$, berlaku $\snorm{f_n - f}_S < \epsilon$.  Karena $\snorm{f_n-f}_S$
merupakan supremum dari $\babs{f_n(x)-f(x)}$, kita melihat bahwa untuk semua $x \in S$,
berlaku $\babs{f_n(x)-f(x)} \leq \snorm{f_n - f}_S < \epsilon$.

Di sisi lain, andaikan $\{ f_n \}_{n=1}^\infty$ konvergen seragam ke $f$.
Misalkan diberikan $\epsilon > 0$.  Pilih $N$ sedemikian sehingga
untuk semua $n \geq N$, berlaku
$\babs{f_n(x)-f(x)} < \epsilon$ untuk semua $x \in S$.
Dengan mengambil supremum atas $x \in S$, kita melihat bahwa
$\snorm{f_n - f}_S \leq \epsilon$.  Jadi, $\lim_{n\to\infty} \snorm{f_n-f}_S = 0$.
\end{proof}

Terkadang dikatakan bahwa
\emph{$\{ f_n \}_{n=1}^\infty$ konvergen ke $f$ dalam norma seragam}%
\index{konvergen dalam norma seragam}\index{konvergensi norma seragam}
alih-alih \emph{konvergen seragam} jika $\snorm{f_n-f}_S \to 0$.  Proposisi tersebut
menyatakan bahwa kedua pengertian itu sama untuk fungsi terbatas.

\begin{example}
Misalkan $f_n \colon [0,1] \to \R$ didefinisikan oleh $f_n(x) \coloneqq \frac{nx+ \sin(nx^2)}{n}$.
Kita menyatakan bahwa $\{ f_n \}_{n=1}^\infty$ konvergen seragam ke $f(x) \coloneqq x$.  Mari kita hitung:
\begin{equation*}
\begin{split}
\snorm{f_n-f}_{[0,1]}
& =
\sup \left\{ \abs{\frac{nx+ \sin(nx^2)}{n} - x} : x \in [0,1] \right\}
\\
& =
\sup \left\{ \frac{\babs{\sin(nx^2)}}{n} : x \in [0,1] \right\}
\\
& \leq
\sup \bigl\{ \nicefrac{1}{n} : x \in [0,1] \bigr\}
\\
& = \nicefrac{1}{n}.
\end{split}
\end{equation*}
\end{example}

Dengan menggunakan norma seragam, kita mendefinisikan barisan Cauchy secara serupa dengan
cara kita mendefinisikan barisan Cauchy bilangan real.

\begin{defn}
Misalkan $f_n \colon S \to \R$ merupakan fungsi terbatas.
Barisan tersebut \emph{\myindex{Cauchy dalam norma seragam}}
atau \emph{\myindex{Cauchy seragam}}
jika untuk setiap $\epsilon > 0$, terdapat $N \in \N$ sedemikian
sehingga untuk semua $m,k \geq N$,
\begin{equation*}
\snorm{f_m-f_k}_S < \epsilon .
\end{equation*}
\end{defn}

\begin{prop} \label{prop:uniformcauchy}
Misalkan $f_n \colon S \to \R$ merupakan fungsi terbatas.
Maka $\{ f_n \}_{n=1}^\infty$ bersifat Cauchy dalam norma seragam jika dan hanya jika
terdapat $f \colon S \to \R$ dan $\{ f_n \}_{n=1}^\infty$ konvergen
seragam ke $f$.
\end{prop}

\begin{proof}
Pertama, andaikan $\{ f_n \}_{n=1}^\infty$ bersifat Cauchy dalam norma seragam.
Mari kita definisikan $f$.  Tetapkan $x$.
Barisan $\bigl\{ f_n(x) \bigr\}_{n=1}^\infty$ bersifat Cauchy karena
\begin{equation*}
\babs{f_m(x)-f_k(x)}
\leq
\snorm{f_m-f_k}_S .
\end{equation*}
Dengan demikian, $\bigl\{ f_n(x) \bigr\}_{n=1}^\infty$ konvergen ke suatu bilangan real.  Definisikan $f \colon S
\to \R$ dengan
\begin{equation*}
f(x) \coloneqq \lim_{n \to \infty} f_n(x) .
\end{equation*}
Barisan
$\{ f_n \}_{n=1}^\infty$ konvergen titik demi titik ke $f$.  Untuk menunjukkan bahwa konvergensinya
seragam, misalkan diberikan $\epsilon > 0$.  Pilih $N$ sedemikian sehingga
untuk semua $m, k \geq N$, berlaku
$\snorm{f_m-f_k}_S < \nicefrac{\epsilon}{2}$.  Dengan kata lain, untuk
semua $x$, berlaku
$\babs{f_m(x)-f_k(x)} < \nicefrac{\epsilon}{2}$.  Untuk sebarang $x$ yang tetap, ambil limit
ketika $k$ menuju tak hingga.  Maka $\babs{f_m(x)-f_k(x)}$
menuju $\babs{f_m(x)-f(x)}$.
Akibatnya, untuk semua $x$,
\begin{equation*}
\babs{f_m(x)-f(x)} \leq \nicefrac{\epsilon}{2} < \epsilon .
\end{equation*}
Jadi, $\{ f_n \}_{n=1}^\infty$ konvergen seragam.

Selanjutnya, kita buktikan arah sebaliknya.
Andaikan $\{ f_n \}_{n=1}^\infty$ konvergen seragam ke
$f$.  Diberikan $\epsilon > 0$, pilih $N$ sedemikian sehingga untuk semua $n \geq N$,
berlaku $\babs{f_n(x)-f(x)} < \nicefrac{\epsilon}{4}$ untuk semua $x \in S$.
Oleh karena itu, untuk semua $m, k \geq N$ dan semua $x$,
\begin{multline*}
\babs{f_m(x)-f_k(x)} = 
\babs{f_m(x)-f(x)+f(x)-f_k(x)}
\\
\leq
\babs{f_m(x)-f(x)}+\babs{f(x)-f_k(x)} < \nicefrac{\epsilon}{4} +
\nicefrac{\epsilon}{4} = \nicefrac{\epsilon}{2} .
\end{multline*}
Ambil supremum atas semua $x$ untuk memperoleh
\begin{equation*}
\snorm{f_m-f_k}_S \leq \nicefrac{\epsilon}{2} < \epsilon .  \qedhere
\end{equation*}
\end{proof}

\subsection{Latihan}

\begin{exercise}
Misalkan $f$ dan $g$ merupakan fungsi terbatas pada $[a,b]$.  Buktikan bahwa
\begin{equation*}
\snorm{f+g}_{[a,b]} \leq \snorm{f}_{[a,b]} + \snorm{g}_{[a,b]} .
\end{equation*}
\end{exercise}

\begin{exercise}
\leavevmode
\begin{enumerate}[a)]
\item
Tentukan limit titik demi titik dari
$\left\{ \dfrac{e^{x/n}}{n} \right\}_{n=1}^\infty$ untuk $x \in \R$.
\item
Apakah limitnya seragam pada $\R$?
\item
Apakah limitnya seragam pada $[0,1]$?
\end{enumerate}
\end{exercise}

\begin{exercise}
Andaikan $f_n \colon S \to \R$ merupakan fungsi yang konvergen seragam
ke $f \colon S \to \R$.  Andaikan $A \subset S$.  Tunjukkan bahwa
barisan pembatasan $\{ f_n|_A \}_{n=1}^\infty$ konvergen seragam ke $f|_A$.
\end{exercise}

\begin{exercise}
Andaikan $\{ f_n \}_{n=1}^\infty$ dan $\{ g_n \}_{n=1}^\infty$ yang terdefinisi pada suatu himpunan $A$ masing-masing konvergen
titik demi titik ke $f$ dan $g$.  Tunjukkan bahwa $\{ f_n+g_n \}_{n=1}^\infty$ konvergen
titik demi titik ke $f+g$.
\end{exercise}

\begin{exercise}
Andaikan $\{ f_n \}_{n=1}^\infty$ dan $\{ g_n \}_{n=1}^\infty$ yang terdefinisi pada suatu himpunan $A$ masing-masing konvergen
seragam ke $f$ dan $g$ pada $A$.
Tunjukkan bahwa $\{ f_n+g_n \}_{n=1}^\infty$
konvergen seragam ke $f+g$ pada $A$.
\end{exercise}

\begin{exercise}
Temukan contoh barisan fungsi $\{ f_n \}_{n=1}^\infty$ dan
$\{ g_n \}_{n=1}^\infty$
yang masing-masing konvergen seragam ke $f$ dan $g$ pada suatu himpunan $A$, tetapi
$\{ f_ng_n \}_{n=1}^\infty$ (hasil kalinya) tidak konvergen seragam ke $fg$ pada $A$.
Petunjuk: Tetapkan $A \coloneqq \R$ dan $f(x)\coloneqq g(x) \coloneqq x$.  Anda bahkan dapat memilih $f_n = g_n$.
\end{exercise}

\begin{exercise}
Andaikan terdapat barisan fungsi $\{ g_n \}_{n=1}^\infty$ yang konvergen
seragam ke $0$ pada $A$.  Sekarang andaikan terdapat barisan fungsi
$\{ f_n \}_{n=1}^\infty$ dan fungsi $f$ pada $A$ sedemikian sehingga
\begin{equation*}
\babs{f_n(x) - f(x)} \leq g_n(x) 
\end{equation*}
untuk semua $x \in A$.  Tunjukkan bahwa $\{ f_n \}_{n=1}^\infty$ konvergen seragam ke $f$ pada $A$.
\end{exercise}

\begin{exercise}
Misalkan $\{ f_n \}_{n=1}^\infty$, $\{ g_n \}_{n=1}^\infty$, dan
$\{ h_n \}_{n=1}^\infty$ merupakan barisan fungsi pada
$[a,b]$.  Andaikan $\{ f_n \}_{n=1}^\infty$ dan $\{ h_n \}_{n=1}^\infty$ konvergen seragam ke suatu fungsi
$f \colon [a,b] \to \R$, serta andaikan $f_n(x) \leq g_n(x) \leq h_n(x)$
untuk semua $x \in [a,b]$.  Tunjukkan bahwa $\{ g_n \}_{n=1}^\infty$ konvergen seragam ke $f$.
\end{exercise}

\begin{exercise}
Misalkan $f_n \colon [0,1] \to \R$ merupakan barisan fungsi naik (artinya,
$f_n(x) \geq f_n(y)$ setiap kali $x \geq y$).  Andaikan $f_n(0) = 0$
dan $\lim\limits_{n \to \infty} f_n(1) = 0$.  Tunjukkan bahwa
$\{ f_n \}_{n=1}^\infty$
konvergen seragam ke $0$.
\end{exercise}

\begin{exercise}
Perhatikan suatu barisan fungsi $f_n \colon [0,1] \to \R$ sedemikian sehingga
terdapat barisan bilangan-bilangan yang saling berbeda $x_n \in [0,1]$ yang memenuhi,
untuk semua $n$,
\begin{equation*}
f_n(x_n) = 1.
\end{equation*}
Buktikan atau sangkal pernyataan-pernyataan berikut:
\begin{enumerate}[a)]
\item
Benar atau salah: Terdapat $\{ f_n \}_{n=1}^\infty$ seperti di atas yang
konvergen titik demi titik ke $0$.
\item
Benar atau salah: Terdapat $\{ f_n \}_{n=1}^\infty$ seperti di atas yang
konvergen seragam ke $0$.
\end{enumerate}
\end{exercise}

\begin{exercise}
Tetapkan fungsi kontinu $h \colon [a,b] \to \R$.
Tetapkan $f(x) \coloneqq h(x)$ untuk $x \in [a,b]$,
$f(x) \coloneqq h(a)$ untuk $x < a$, dan $f(x) \coloneqq h(b)$ untuk semua $x > b$.  Pertama, tunjukkan
bahwa $f \colon \R \to \R$ kontinu.
Sekarang misalkan $f_n$ adalah
fungsi $g$ dari \exerciseref{exercise:smoothingout} dengan
$\epsilon = \nicefrac{1}{n}$, yang didefinisikan pada interval $[a,b]$.  Artinya,
\begin{equation*}
f_n(x) \coloneqq \frac{n}{2} \int_{x-1/n}^{x+1/n} f .
\end{equation*}
Tunjukkan bahwa $\{ f_n \}_{n=1}^\infty$ konvergen seragam ke $h$ pada $[a,b]$.
\end{exercise}


\begin{exercise}
Buktikan bahwa
jika barisan fungsi $f_n \colon S \to \R$
konvergen seragam ke fungsi terbatas $f \colon S \to \R$,
maka terdapat $N$ sedemikian sehingga untuk semua $n \geq N$, fungsi $f_n$
terbatas.
\end{exercise}

\begin{exercise}
Andaikan terdapat satu konstanta $B$ dan
barisan fungsi
$f_n \colon S \to \R$ yang dibatasi oleh $B$,
yaitu $\babs{f_n(x)} \leq B$ untuk semua $x \in S$.
Andaikan $\{ f_n \}_{n=1}^\infty$ konvergen titik demi titik
ke $f \colon S \to \R$.
Buktikan bahwa $f$ terbatas.
\end{exercise}

\begin{exercise}[memerlukan \sectionref{sec:moreonseries}]
Dalam \exampleref{example:geomsumptconv}, kita melihat bahwa
$\sum_{k=0}^\infty x^k$ konvergen titik demi titik ke $\frac{1}{1-x}$ pada
$(-1,1)$.
\begin{enumerate}[a)]
\item
Tunjukkan bahwa untuk setiap $c$ dengan $0 \leq c < 1$, deret
$\sum_{k=0}^\infty x^k$ konvergen seragam pada $[-c,c]$.
\item
Tunjukkan bahwa deret $\sum_{k=0}^\infty x^k$ tidak konvergen seragam
pada $(-1,1)$.
\end{enumerate}
\end{exercise}

%%%%%%%%%%%%%%%%%%%%%%%%%%%%%%%%%%%%%%%%%%%%%%%%%%%%%%%%%%%%%%%%%%%%%%%%%%%%%%

\sectionnewpage
\section{Pertukaran limit}
\label{sec:liminter}

%mbxINTROSUBSECTION

\sectionnotes{1--2.5 kuliah;
subbagian tentang turunan dan deret pangkat (yang
memerlukan \sectionref{sec:moreonseries}) bersifat opsional.}

Sebagian besar analisis modern terutama membahas persoalan
pertukaran dua operasi limit.  Ketika
kita memiliki rangkaian dua limit, kita tidak selalu dapat langsung menukar keduanya.
Sebagai contoh,
\begin{equation*}
0 = 
\lim_{n\to\infty}
\left(
\lim_{k\to\infty}
\frac{n}{n + k}
\right)
\neq
\lim_{k\to\infty}
\left(
\lim_{n\to\infty}
\frac{n}{n + k}
\right)
= 1 .
\end{equation*}

Ketika membahas barisan fungsi, pertukaran limit cukup sering muncul.
Kita akan melihat beberapa kasus: kekontinuan limit,
integral limit, turunan limit, dan konvergensi
deret pangkat.

\subsection{Kekontinuan limit}

Jika kita memiliki barisan $\{ f_n \}_{n=1}^\infty$ yang terdiri atas fungsi kontinu, apakah limitnya kontinu?
Andaikan $f$ merupakan limit (titik demi titik) dari $\{ f_n \}_{n=1}^\infty$.
Jika $\lim_{k\to\infty} x_k = x$,
kita tertarik pada
pertukaran limit berikut, dengan kesamaan yang hendak dibuktikan
ditandai oleh tanda tanya.  Kesamaan tersebut tidak selalu benar.
Bahkan, limit-limit di sebelah kiri
tanda tanya mungkin tidak ada.
\begin{equation*}
\lim_{k \to \infty} 
f(x_k)
=
\lim_{k \to \infty} 
\Bigl(
\lim_{n \to \infty} f_n(x_k)
\Bigr)
\overset{?}{=}
\lim_{n \to \infty}
\Bigl(
\lim_{k \to \infty} 
f_n(x_k)
\Bigr)
=
\lim_{n \to \infty}
f_n(x)
=
f(x) .
\end{equation*}
Kita ingin menemukan syarat pada barisan $\{ f_n \}_{n=1}^\infty$
agar persamaan di atas berlaku.
Jika kita hanya mensyaratkan konvergensi titik demi titik, maka limit
suatu barisan fungsi belum tentu kontinu dan persamaan di atas
belum tentu berlaku.

\begin{example}
Definisikan $f_n \colon [0,1] \to \R$ dengan
\begin{equation*}
f_n(x) \coloneqq
\begin{cases}
1-nx &  \text{jika } x < \nicefrac{1}{n},\\
0 &  \text{jika } x \geq \nicefrac{1}{n}.
\end{cases}
\end{equation*}
Lihat \figureref{contconvcntr:fig}.

\begin{myfigureht}
\myincludegraphics{contconvcntr}{%
Grafik suatu fungsi pada interval
dari 0 hingga 1; fungsi tersebut bernilai 1 di titik ujung kiri, kemudian
grafiknya berupa garis lurus yang turun ke 0 pada sumbu x di titik x sama dengan satu per n.
Setelah itu, fungsi bernilai nol hingga titik ujung kanan.}
\caption{Grafik $f_n(x)$.%
\label{contconvcntr:fig}}
\end{myfigureht}

Setiap fungsi $f_n$ kontinu.
Tetapkan $x \in (0,1]$.  Jika $n \geq \nicefrac{1}{x}$,
maka $x \geq \nicefrac{1}{n}$.  Oleh karena itu, untuk $n \geq \nicefrac{1}{x}$,
berlaku $f_n(x) = 0$, sehingga
\begin{equation*}
\lim_{n \to \infty} f_n(x) = 0.
\end{equation*}
Di sisi lain, jika $x=0$, maka
\begin{equation*}
\lim_{n \to \infty} f_n(0) = 
\lim_{n \to \infty} 1 = 1.
\end{equation*}
Dengan demikian, limit titik demi titik dari $f_n$ adalah fungsi
$f \colon [0,1] \to \R$ yang didefinisikan oleh
\begin{equation*}
f(x) \coloneqq
\begin{cases}
1 &  \text{jika } x = 0,\\
0 &  \text{jika } x > 0.
\end{cases}
\end{equation*}
Fungsi $f$ tidak kontinu di 0.
\end{example}

Namun, jika kita mensyaratkan konvergensinya seragam, limit-limit tersebut dapat
dipertukarkan.

\begin{thm}
Andaikan $S \subset \R$.
Misalkan $\{ f_n \}_{n=1}^\infty$ merupakan
barisan fungsi kontinu $f_n \colon S \to \R$ yang konvergen
seragam ke $f \colon S \to \R$.  Maka $f$ kontinu.
\end{thm}

\begin{proof}
Tetapkan $x \in S$.  Misalkan $\{ x_n \}_{n=1}^\infty$ merupakan barisan dalam $S$
yang konvergen ke $x$.

Misalkan diberikan $\epsilon > 0$.
Karena $\{ f_k \}_{k=1}^\infty$ konvergen seragam ke $f$, kita memperoleh $k \in \N$ sedemikian sehingga
\begin{equation*}
\babs{f_k(y)-f(y)} < \nicefrac{\epsilon}{3}
\end{equation*}
untuk semua $y \in S$.  Karena $f_k$ kontinu di $x$,
kita memperoleh $N \in \N$ sedemikian sehingga untuk semua $m \geq N$,
\begin{equation*}
\babs{f_k(x_m)-f_k(x)} < \nicefrac{\epsilon}{3} .
\end{equation*}
Dengan demikian, untuk semua
$m \geq N$,
\begin{equation*}
\begin{split}
\babs{f(x_m)-f(x)}
& =
\babs{f(x_m)-f_k(x_m)+f_k(x_m)-f_k(x)+f_k(x)-f(x)}
\\
& \leq
\babs{f(x_m)-f_k(x_m)}+
\babs{f_k(x_m)-f_k(x)}+
\babs{f_k(x)-f(x)}
\\
& <
\nicefrac{\epsilon}{3} +
\nicefrac{\epsilon}{3} +
\nicefrac{\epsilon}{3} = \epsilon .
\end{split}
\end{equation*}
Oleh karena itu, $\bigl\{ f(x_m) \bigr\}_{m=1}^\infty$ konvergen ke $f(x)$,
dan akibatnya $f$ kontinu di $x$.
Karena $x$ sebarang, $f$ kontinu di setiap titik.
\end{proof}

\subsection{Integral limit}

Sekali lagi, jika kita hanya mensyaratkan konvergensi titik demi titik, maka integral
limit suatu barisan fungsi belum tentu sama dengan limit
integral-integralnya.

\begin{example}
Definisikan $f_n \colon [0,1] \to \R$ dengan
\begin{equation*}
f_n(x) \coloneqq
\begin{cases}
0 &  \text{jika } x = 0,\\
n-n^2x &  \text{jika } 0 < x < \nicefrac{1}{n},\\
0 &  \text{jika } x \geq \nicefrac{1}{n}.
\end{cases}
\avoidbreak
\end{equation*}
Lihat \figureref{intconvcntr:fig}.

\begin{myfigureht}
\myincludegraphics{intconvcntr}{%
Grafik suatu fungsi pada interval 0 hingga 1 yang bernilai 0 di titik ujung kiri,
tetapi kemudian memiliki diskontinuitas loncat dan dinyatakan oleh
garis menurun dari titik (0,n) ke sumbu x
di titik x sama dengan satu per n.  Setelah itu, fungsi bernilai 0 pada sisa interval.}
\caption{Grafik $f_n(x)$.%
\label{intconvcntr:fig}}
\end{myfigureht}

Setiap $f_n$ terintegralkan secara Riemann (fungsi tersebut kontinu pada $(0,1]$ dan terbatas),
dan teorema dasar kalkulus menyatakan bahwa
\begin{equation*}
\int_0^1 f_n =
\int_0^{\nicefrac{1}{n}} (n-n^2x)\,dx = \nicefrac{1}{2} .
\end{equation*}
Mari kita hitung limit titik demi titik dari $\{ f_n \}_{n=1}^\infty$.
Tetapkan $x \in (0,1]$.  Untuk $n \geq \nicefrac{1}{x}$,
berlaku $x \geq \nicefrac{1}{n}$ dan karena itu $f_n(x) = 0$.  Jadi,
\begin{equation*}
\lim_{n \to \infty} f_n(x) = 0.
\end{equation*}
Kita juga memiliki $f_n(0) = 0$ untuk semua $n$.  Jadi, limit titik demi titik
dari $\{ f_n \}_{n=1}^\infty$ adalah fungsi nol.
Ringkasnya,
\begin{equation*}
\nicefrac{1}{2} =
\lim_{n\to\infty}
\int_0^1 f_n (x)\,dx
\neq
\int_0^1
\left(
\lim_{n\to\infty}
f_n(x)\right)\,dx
=
\int_0^1 0\,dx = 0 .
\end{equation*}
\end{example}

Namun, jika kita mensyaratkan konvergensinya seragam, limit-limit tersebut dapat
dipertukarkan.\footnote{Syarat yang lebih lemah
sudah cukup untuk teorema semacam ini, tetapi pembuktian generalisasi tersebut memerlukan
perangkat yang lebih canggih daripada yang dibahas di sini: integral Lebesgue.
Khususnya, teorema ini berlaku dengan konvergensi titik demi titik
asalkan $f$ terintegralkan dan terdapat $M$ sedemikian sehingga
$\snorm{f_n}_{[a,b]} \leq M$ untuk semua $n$.}

\begin{thm} \label{integralinterchange:thm}
Misalkan $\{ f_n \}_{n=1}^\infty$ merupakan barisan fungsi yang terintegralkan secara Riemann,
$f_n \colon [a,b] \to \R$
yang konvergen seragam ke $f \colon [a,b]
\to \R$.  Maka $f$ terintegralkan secara Riemann dan
\begin{equation*}
\int_a^b f = \lim_{n\to\infty} \int_a^b f_n .
\end{equation*}
\end{thm}

\begin{proof}
Misalkan diberikan $\epsilon > 0$.
Karena $f_n$ menuju $f$ secara seragam, kita memperoleh $M \in \N$ sedemikian sehingga
untuk semua $n \geq M$, berlaku
$\babs{f_n(x)-f(x)} < \frac{\epsilon}{2(b-a)}$ untuk semua $x \in [a,b]$.
Khususnya, berdasarkan pertidaksamaan segitiga terbalik,
$\babs{f(x)} < \frac{\epsilon}{2(b-a)} + \babs{f_n(x)}$ untuk semua $x$.
Jadi, $f$ terbatas,
karena $f_n$ terbatas.
Perhatikan bahwa $f_n$ terintegralkan, lalu hitung
\begin{equation*}
\begin{split}
\overline{\int_a^b} f
-
\underline{\int_a^b} f
& =
\overline{\int_a^b} \bigl( f(x) - f_n(x) + f_n(x) \bigr)\,dx
-
\underline{\int_a^b} \bigl( f(x) - f_n(x) + f_n(x) \bigr)\,dx
\\
& \leq
\overline{\int_a^b} \bigl( f(x) - f_n(x) \bigr)\,dx +  \overline{\int_a^b} f_n(x) \,dx
-
\underline{\int_a^b} \bigl( f(x) - f_n(x) \bigr)\,dx -  \underline{\int_a^b} f_n(x) \,dx
\\
& =
\overline{\int_a^b} \bigl( f(x) - f_n(x) \bigr)\,dx +  \int_a^b f_n(x) \,dx
-
\underline{\int_a^b} \bigl( f(x) - f_n(x) \bigr)\,dx -  \int_a^b f_n(x) \,dx
\\
& =
\overline{\int_a^b} \bigl( f(x) - f_n(x) \bigr)\,dx
-
\underline{\int_a^b} \bigl( f(x) - f_n(x) \bigr)\,dx
\\
& \leq
\frac{\epsilon}{2(b-a)} (b-a) + 
\frac{\epsilon}{2(b-a)} (b-a) = \epsilon .
\end{split}
\end{equation*}
Pertidaksamaan pertama mengikuti dari \propref{prop:upperlowerlinineq}.
Pertidaksamaan kedua mengikuti \propref{intulbound:prop} dan
fakta bahwa untuk semua $x \in [a,b]$, berlaku
$\frac{-\epsilon}{2(b-a)} < f(x)-f_n(x) < \frac{\epsilon}{2(b-a)}$.
Karena $\epsilon > 0$ sebarang, $f$ terintegralkan secara Riemann.

Terakhir, kita menghitung $\int_a^b f$.  Kita menerapkan \propref{intbound:prop}
dalam perhitungan.  Sekali lagi, untuk semua $n \geq M$ (dengan $M$ sama seperti di atas),
\begin{equation*}
\begin{split}
\abs{\int_a^b f - \int_a^b f_n} & = 
\abs{ \int_a^b \bigl(f(x) - f_n(x)\bigr)\,dx}
\\
& \leq
\frac{\epsilon}{2(b-a)} (b-a) = \frac{\epsilon}{2} < \epsilon .
\end{split}
\end{equation*}
Oleh karena itu, $\bigl\{ \int_a^b f_n \bigr\}_{n=1}^\infty$ konvergen ke $\int_a^b f$.
\end{proof}

\begin{example}
Andaikan kita ingin menghitung
\begin{equation*}
\lim_{n\to\infty} \int_0^1 \frac{nx+ \sin(nx^2)}{n} \,dx .
\end{equation*}
Integral untuk suatu $n$ tertentu mustahil dihitung dengan
kalkulus karena $\sin(nx^2)$ tidak memiliki antiturunan dalam bentuk tertutup.  Namun,
kita dapat menghitung limitnya.
Kita telah menunjukkan sebelumnya bahwa $\frac{nx+ \sin(nx^2)}{n}$ konvergen seragam
pada $[0,1]$ ke $x$.
Berdasarkan \thmref{integralinterchange:thm}, limit tersebut ada dan
\begin{equation*}
\lim_{n\to\infty} \int_0^1 \frac{nx+ \sin(nx^2)}{n} \,dx
=
\int_0^1
x \,dx = \nicefrac{1}{2} .
\end{equation*}
\end{example}

\begin{example}
Jika konvergensinya hanya titik demi titik, limitnya bahkan belum tentu terintegralkan secara
Riemann.  Pada $[0,1]$, definisikan
\begin{equation*}
f_n(x) \coloneqq
\begin{cases}
1 & \text{jika } x = \nicefrac{p}{q} \text{ dalam bentuk paling sederhana dan } q \leq n, \\
0 & \text{selain itu.}
\end{cases}
\end{equation*}
Setiap fungsi $f_n$ berbeda dari fungsi nol hanya pada berhingga banyak titik;
hanya ada berhingga banyak pecahan dalam $[0,1]$ dengan penyebut kurang dari
atau sama dengan $n$.   Jadi, $f_n$ terintegralkan dan $\int_0^1 f_n = \int_0^1 0 =
0$.  Sebagai latihan sederhana, tunjukkan bahwa $\{ f_n \}_{n=1}^\infty$ konvergen titik demi titik ke
\myindex{fungsi Dirichlet}
\begin{equation*}
f(x) \coloneqq
\begin{cases}
1 & \text{jika } x \in \Q, \\
0 & \text{selain itu,}
\end{cases}
\end{equation*}
yang tidak terintegralkan secara Riemann.
\end{example}

\begin{example}
Bahkan, jika konvergensinya hanya titik demi titik, limit fungsi-fungsi
terbatas belum tentu terbatas.
Definisikan $f_n \colon [0,1] \to \R$ dengan
\begin{equation*}
f_n(x) \coloneqq
\begin{cases}
0 & \text{jika } x < \nicefrac{1}{n},\\
\nicefrac{1}{x} & \text{selain itu.}
\end{cases}
\end{equation*}
Untuk setiap $n$, berlaku $\babs{f_n(x)} \leq n$ untuk semua $x \in [0,1]$, sehingga
fungsi-fungsi tersebut terbatas.  Namun, $\{ f_n \}_{n=1}^\infty$ konvergen titik demi titik ke
fungsi tak terbatas
\begin{equation*}
f(x) \coloneqq
\begin{cases}
0 & \text{jika } x = 0,\\
\nicefrac{1}{x} & \text{selain itu.}
\end{cases}
\end{equation*}
\end{example}

\subsection{Turunan limit}

Walaupun konvergensi seragam cukup untuk menukar limit dengan integral, konvergensi tersebut tidak
cukup untuk menukar limit dengan turunan, kecuali jika turunan-turunannya sendiri juga
konvergen seragam.

\begin{example}
Misalkan $f_n(x) \coloneqq \frac{\sin(nx)}{n}$.  Maka $f_n$ konvergen seragam ke
0.  Lihat \figureref{fig:conv1nsinxn}.
Turunan limitnya adalah 0.  Namun, $f_n'(x) = \cos(nx)$, yang
bahkan tidak konvergen titik demi titik.
Sebagai contoh, $f_n'(\pi) = {(-1)}^n$.
Selain itu, $f_n'(0) = 1$ untuk semua $n$, yang memang konvergen, tetapi tidak ke $0$.
\begin{myfigureht}
\myincludegraphics{conv1nsinxn}{%
Gelombang sinus untuk 10 frekuensi berbeda, tetapi kini amplitudonya juga diperkecil
untuk frekuensi yang lebih tinggi.  Fungsi berfrekuensi lebih tinggi
lebih dekat ke sumbu x, sehingga semua grafik itu bersama-sama tidak lagi memenuhi seluruh pita.}
\caption{Grafik $\frac{\sin(nx)}{n}$ untuk
$n=1,2,\ldots,10$, dengan $n$ yang lebih besar digambar dalam warna abu-abu yang lebih muda.%
\label{fig:conv1nsinxn}}
\end{myfigureht}
\end{example}

\begin{example} \label{exercise:badconvergenceder}
Misalkan $f_n(x) \coloneqq \frac{1}{1+nx^2}$.
Jika $x \neq 0$, maka $\lim_{n \to \infty} f_n(x) = 0$,
tetapi $\lim_{n \to \infty} f_n(0) = 1$.
Oleh karena itu, $\{ f_n \}_{n=1}^\infty$ konvergen titik demi titik menuju fungsi yang tidak kontinu
di $0$.
Kita hitung
\begin{equation*}
f_n'(x) %= \frac{d}{dx} \frac{1}{1+n x^2}
= \frac{-2 n x}{(1+ n x^2)^2} .
\end{equation*}
Untuk setiap $x$, $\lim_{n\to\infty} f_n'(x) = 0$, sehingga turunan-turunannya
konvergen titik demi titik ke 0,
tetapi pembaca dapat memeriksa bahwa konvergensinya tidak seragam pada interval mana pun
yang memuat $0$.
Limit $f_n$ tidak terdiferensialkan di $0$---bahkan tidak
kontinu di $0$.  Lihat \figureref{fig:convder1over1pnxsq}.
\begin{myfigureht}
\myincludegraphics{convder1over1pnxsq}{%
Dua plot, masing-masing memperlihatkan beberapa fungsi
dalam berbagai corak abu-abu.
Di sebelah kiri terdapat plot kurva berbentuk lonceng, dengan puncak yang semakin sempit
(tetapi tetap sama tinggi) untuk corak abu-abu yang lebih muda.
Di sebelah kanan terdapat turunannya; grafik bernilai positif dengan puncak sebelum
titik asal, lalu melewati titik asal dan mencapai minimum setelah
titik asal sebelum kembali naik menuju sumbu x.  Semakin muda
warnanya, puncak dan minimumnya semakin sempit dan semakin jauh
dari sumbu x.}
\caption{Grafik $\frac{1}{1+nx^2}$ dan turunannya
untuk $n=1,2,\ldots,10$, dengan $n$ yang lebih besar digambar dalam warna abu-abu yang lebih muda.%
\label{fig:convder1over1pnxsq}}
\end{myfigureht}
\end{example}

Lihat latihan untuk contoh-contoh lainnya.  Dengan menggunakan
teorema dasar kalkulus, kita memperoleh jawaban untuk fungsi-fungsi yang
terdiferensialkan secara kontinu.  Teorema berikut tetap benar sekalipun
kita tidak mengandaikan bahwa turunan-turunannya kontinu, tetapi pembuktiannya lebih
sulit.

\begin{thm} \label{thm:dersconverge}
Misalkan $I$ adalah interval terbatas dan
$f_n \colon I \to \R$ merupakan fungsi-fungsi yang terdiferensialkan secara kontinu.
Andaikan $\{ f_n' \}_{n=1}^\infty$ konvergen seragam ke $g \colon I \to \R$,
dan andaikan $\bigl\{ f_n(c) \bigr\}_{n=1}^\infty$ merupakan
barisan konvergen untuk suatu $c \in I$.  Maka $\{ f_n \}_{n=1}^\infty$ konvergen seragam menuju
fungsi $f \colon I \to \R$ yang terdiferensialkan secara kontinu, dan $f' = g$.
\end{thm}

\begin{proof}
Definisikan $f(c) \coloneqq \lim_{n\to \infty} f_n(c)$.
Karena $f_n'$ kontinu dan karenanya terintegralkan Riemann,
dengan teorema dasar kalkulus kita peroleh, untuk $x \in I$,
\begin{equation*}
f_n(x) = f_n(c) + \int_c^x f_n' .
\end{equation*}
Karena $\{ f_n' \}_{n=1}^\infty$ konvergen seragam pada $I$, barisan tersebut konvergen seragam
pada $[c,x]$ (atau $[x,c]$ jika $x < c$).
Dengan demikian, limit ketika $n \to \infty$ pada ruas kanan ada.
Definisikan $f$ pada titik-titik lainnya (dengan $x\neq c$) melalui limit ini:
\begin{equation*}
f(x) \coloneqq
\lim_{n\to\infty} f_n(c) + \lim_{n\to\infty} \int_c^x f_n'
=
f(c) + \int_c^x g .
\end{equation*}
Fungsi $g$ kontinu karena merupakan limit seragam dari fungsi-fungsi
kontinu.  Oleh karena itu, $f$ terdiferensialkan dan $f'(x) = g(x)$ untuk semua $x \in I$
berdasarkan bentuk kedua teorema dasar kalkulus.

Tinggal membuktikan
konvergensi seragam.
Andaikan $I$ memiliki batas bawah $a$ dan batas atas $b$.
Ambil sebarang $\epsilon > 0$.  Pilih $M$
sedemikian sehingga untuk semua $n \geq M$ berlaku
$\babs{f(c)-f_n(c)} < \nicefrac{\epsilon}{2}$
dan
$\babs{g(x)-f_n'(x)} < \frac{\epsilon}{2(b-a)}$
untuk semua $x \in I$.  Maka
\begin{equation*}
\begin{split}
\babs{f(x) - f_n(x)} & =
\abs{\left(f(c) + \int_c^x g\right) - \left( f_n(c) + \int_c^x f_n' \right)}
\\
& \leq
\babs{f(c) - f_n(c)} + \abs{\int_c^x g - \int_c^x f_n'}
\\
& =
\babs{f(c) - f_n(c)} + \abs{\int_c^x \bigl(g(s) - f_n'(s)\bigr) \, ds}
\\
& <
\frac{\epsilon}{2}
+
\frac{\epsilon}{2(b-a)}
(b-a)
=\epsilon. \qedhere
\end{split}
\end{equation*}
\end{proof}

Pembuktian tersebut tetap berlaku tanpa asumsi bahwa $I$ terbatas, kecuali untuk
konvergensi seragam $f_n$ ke $f$.  Sebagai contoh, andaikan $I = \R$ dan definisikan
$f_n(x) \coloneqq \nicefrac{x}{n}$.  Maka $f_n'(x)=\nicefrac{1}{n}$, yang
konvergen seragam ke $0$.  Namun, $\{f_n\}_{n=1}^\infty$ hanya konvergen titik demi titik ke 0.

\subsection{Konvergensi deret pangkat}

Dalam \sectionref{sec:moreonseries} kita melihat bahwa deret pangkat konvergen
mutlak di dalam jari-jari konvergensinya, sehingga deret tersebut konvergen titik demi titik.
Mari kita tunjukkan bahwa deret itu (beserta semua turunannya) juga konvergen seragam
pada setiap interval tertutup yang termuat di dalam interval konvergensinya.
Fakta ini memungkinkan kita
menukar beberapa jenis limit.  Limitnya bukan hanya kontinu,
tetapi kita juga dapat melakukan integrasi dan bahkan diferensiasi suku demi suku
terhadap deret pangkat yang konvergen.

\begin{prop}
Misalkan $\sum_{n=0}^\infty c_n {(x-a)}^n$ merupakan deret pangkat konvergen dengan
jari-jari konvergensi $\rho$, dengan $0 < \rho \leq \infty$.
Maka deret tersebut konvergen seragam
pada $[a-r,a+r]$ setiap kali $0 < r < \rho$.

Secara khusus, deret tersebut konvergen (titik demi titik) menuju fungsi kontinu
pada $(a-\rho,a+\rho)$ jika $\rho < \infty$, atau pada $\R$ jika $\rho = \infty$.
\end{prop}

\begin{proof}
Misalkan $I \coloneqq (a-\rho,a+\rho)$ jika $\rho < \infty$,
atau $I \coloneqq \R$ jika $\rho= \infty$.
Ambil $0 < r < \rho$.
Deret tersebut konvergen mutlak untuk setiap $x \in I$,
khususnya untuk $x = a+r$.
Jadi, $\sum_{n=0}^\infty \sabs{c_n} r^n$ konvergen.
Untuk sebarang $\epsilon >0$, pilih $M$ sedemikian sehingga untuk semua $k \geq M$,
\begin{equation*}
\sum_{n=k+1}^\infty \sabs{c_n} {r}^n < \epsilon .
\end{equation*}
Untuk semua $x \in [a-r,a+r]$ dan semua $m > k$,
\begin{multline*}
\abs{\sum_{n=0}^m c_n {(x-a)}^n - 
\sum_{n=0}^k c_n {(x-a)}^n}
=
\abs{\sum_{n=k+1}^m c_n {(x-a)}^n}
\\
\leq
\sum_{n=k+1}^m \sabs{c_n} {\sabs{x-a}}^n
\leq
\sum_{n=k+1}^m \sabs{c_n} {r}^n
\leq
\sum_{n=k+1}^\infty \sabs{c_n} {r}^n
<\epsilon.
\end{multline*}
Oleh karena itu, jumlah-jumlah parsialnya bersifat Cauchy seragam pada $[a-r,a+r]$ dan
karenanya konvergen seragam pada himpunan tersebut.

Selain itu, jumlah-jumlah parsialnya merupakan polinom, yang
kontinu, sehingga limit seragamnya pada $[a-r,a+r]$
merupakan fungsi kontinu.
Karena $r < \rho$ dipilih sebarang, fungsi limitnya
kontinu pada seluruh $I$.
\end{proof}

Seperti telah disebutkan, kita akan menunjukkan bahwa
deret pangkat dapat didiferensialkan dan diintegralkan
suku demi suku.  Deret hasil
diferensiasi atau integrasi itu kembali merupakan deret pangkat,
dan kita akan menunjukkan bahwa jari-jari konvergensinya tetap sama.
Oleh karena itu,
setiap deret pangkat mendefinisikan fungsi yang diferensiabel tak berhingga kali.

Pertama-tama kita buktikan bahwa kita dapat mencari antiturunannya, karena integrasi hanya memerlukan
limit seragam.

\begin{cor}
Misalkan $\sum_{n=0}^\infty c_n {(x-a)}^n$ merupakan deret pangkat konvergen dengan
jari-jari konvergensi $0 < \rho \leq \infty$.
Misalkan $I \coloneqq (a-\rho,a+\rho)$ jika $\rho < \infty$
atau $I \coloneqq \R$ jika $\rho= \infty$.  Misalkan $f \colon I \to \R$ merupakan limitnya.
Maka
\begin{equation*}
\int_a^x f = \sum_{n=1}^\infty \frac{c_{n-1}}{n} {(x-a)}^{n} ,
\end{equation*}
dengan jari-jari konvergensi deret ini sekurang-kurangnya $\rho$.
\end{cor}

\begin{proof}
Ambil $0 < r < \rho$.
Jumlah-jumlah parsial $\sum_{n=0}^k c_n {(x-a)}^n$ konvergen seragam pada $[a-r,a+r]$.
Untuk setiap $x \in [a-r,a+r]$ yang tetap, konvergensinya juga seragam
pada $[a,x]$ (atau $[x,a]$ jika $x < a$).
Oleh karena itu,
\begin{equation*}
\int_a^x f =
\int_a^x \lim_{k\to\infty} \sum_{n=0}^k c_n {(s-a)}^n \, ds
=
\lim_{k\to\infty}
\int_a^x \sum_{n=0}^k c_n {(s-a)}^n \, ds
=
\lim_{k\to\infty}
\sum_{n=1}^{k+1} \frac{c_{n-1}}{n} {(x-a)}^{n} . \qedhere
\end{equation*}
\end{proof}


\begin{cor} \label{cor:differentiatepowerser}
Misalkan $\sum_{n=0}^\infty c_n {(x-a)}^n$ merupakan deret pangkat konvergen dengan
jari-jari konvergensi $0 < \rho \leq \infty$.
Misalkan $I \coloneqq (a-\rho,a+\rho)$ jika $\rho < \infty$
atau $I \coloneqq \R$ jika $\rho= \infty$.  Misalkan $f \colon I \to \R$ merupakan limitnya.
Maka $f$ merupakan fungsi diferensiabel, dan
\begin{equation*}
f'(x) = \sum_{n=0}^\infty (n+1) c_{n+1} {(x-a)}^{n} ,
\end{equation*}
dan jari-jari konvergensi deret ini adalah $\rho$.
\end{cor}

\begin{proof}
Ambil $0 < r < \rho$.
Deret tersebut konvergen seragam pada
$[a-r,a+r]$,
tetapi kita perlu menunjukkan bahwa deret turunannya konvergen seragam.
Misalkan
\begin{equation*}
R \coloneqq \limsup_{n \to \infty} \sabs{c_n}^{1/n} .
\end{equation*}
Karena deret tersebut konvergen, $R < \infty$.
Jari-jari konvergensi $\rho$ adalah $\nicefrac{1}{R}$ (atau $\rho = \infty$ jika $R=0$).

Ambil sebarang $\epsilon > 0$.  Dalam \exampleref{example:nto1overn},
kita melihat bahwa $\lim_{n\to\infty} n^{1/n} = 1$.
Oleh karena itu, terdapat $N$ sedemikian sehingga untuk semua $n \geq N$ berlaku
$n^{1/n} < 1+\epsilon$.
Jadi,
\begin{equation*}
R = 
\limsup_{n \to \infty}
\sabs{c_n}^{1/n}
\leq
\limsup_{n \to \infty}
\sabs{n c_n}^{1/n}
\leq
(1+\epsilon)
\limsup_{n \to \infty}
\sabs{c_n}^{1/n}
=
(1+\epsilon)R .
\end{equation*}
Karena $\epsilon$ dipilih sebarang, $\limsup_{n \to \infty} \sabs{n c_n}^{1/n} = R$.
Oleh karena itu, $\sum_{n=1}^\infty n c_{n} {(x-a)}^{n}$ memiliki jari-jari
konvergensi $\rho$.  Dengan membagi deret tersebut dengan $(x-a)$, kita peroleh bahwa
$\sum_{n=0}^\infty (n+1) c_{n+1} {(x-a)}^{n}$ juga memiliki jari-jari konvergensi
$\rho$.

Akibatnya, jumlah-jumlah parsial
$\sum_{n=0}^k (n+1) c_{n+1} {(x-a)}^{n}$,
yang merupakan turunan dari jumlah-jumlah parsial
$\sum_{n=0}^{k+1} c_{n} {(x-a)}^{n}$,
konvergen seragam pada $[a-r,a+r]$.  Selain itu,
deret tersebut jelas konvergen di $x=a$.
Dengan demikian, kita dapat menerapkan \thmref{thm:dersconverge}, dan
pembuktiannya selesai karena $r < \rho$ dipilih sebarang.
\end{proof}

\begin{example} \label{example:exponentialbypowerseries}
Kita dapat saja menggunakan hasil ini untuk mendefinisikan fungsi eksponensial.  Yakni,
deret pangkat
\begin{equation*}
f(x) \coloneqq \sum_{n=0}^\infty \frac{x^n}{n!}
\end{equation*}
memiliki jari-jari konvergensi $\rho=\infty$.  Selain itu,
$f(0) = 1$.
Dengan mendiferensialkannya
suku demi suku, kita peroleh bahwa $f'(x) = f(x)$.
\end{example}

\begin{example}
Deret
\begin{equation*}
\sum_{n=1}^\infty n x^n
\end{equation*}
konvergen ke $\frac{x}{{(1-x)}^2}$ pada $(-1,1)$.

Bukti:
Pada $(-1,1)$, $\sum_{n=0}^\infty x^n$ konvergen ke
$\frac{1}{1-x}$.  Turunannya, $\sum_{n=1}^\infty n x^{n-1}$, kemudian konvergen
pada interval yang sama ke $\frac{1}{{(1-x)}^2}$.  Dengan mengalikannya dengan $x$,
kita memperoleh hasil tersebut.
\end{example}

\subsection{Latihan}

\begin{exercise}
%While uniform convergence preserves continuity, it does not preserve
%differentiability.
Tentukan contoh eksplisit suatu barisan
fungsi diferensiabel pada $[-1,1]$ yang konvergen seragam menuju
suatu fungsi $f$ sedemikian sehingga $f$ tidak diferensiabel.
Petunjuk:
Ada banyak kemungkinan.
Yang paling sederhana mungkin adalah menggabungkan $\sabs{x}$ dan $\frac{n}{2}x^2 +
\frac{1}{2n}$.
Kemungkinan lainnya adalah
mempertimbangkan $\sqrt{x^2+{(\nicefrac{1}{n})}^2}$.  Tunjukkan bahwa fungsi-fungsi ini diferensiabel,
konvergen seragam, lalu tunjukkan bahwa limitnya tidak diferensiabel.
\end{exercise}

\begin{exercise}
Misalkan $f_n(x) \coloneqq \frac{x^n}{n}$.  Tunjukkan bahwa $\{ f_n \}_{n=1}^\infty$ konvergen seragam menuju
fungsi diferensiabel $f$ pada $[0,1]$ (tentukan $f$).  Namun, tunjukkan bahwa
$f'(1) \neq \lim\limits_{n\to\infty} f_n'(1)$.
\end{exercise}

%\begin{exnote}
%Note: The previous two exercises show that
%we cannot simply swap limits with derivatives, even if the convergence is
%uniform.  See also \exerciseref{c1uniflim:exercise} below.
%\end{exnote}
%
\begin{exercise}
Misalkan $f \colon [0,1] \to \R$ terintegralkan Riemann (dan karenanya terbatas).
Tentukan $\displaystyle \lim_{n\to\infty} \int_0^1 \frac{f(x)}{n} \,dx$.
\end{exercise}

\begin{exercise}
Tunjukkan bahwa
$\displaystyle \lim_{n\to\infty} \int_1^2 e^{-nx^2} \,dx = 0$.  Anda boleh
menggunakan fakta-fakta kalkulus tentang fungsi eksponensial.
%what you know about the exponential function from calculus.
\end{exercise}

\begin{exercise}
Tentukan contoh suatu barisan fungsi kontinu pada $(0,1)$ yang konvergen
titik demi titik menuju fungsi kontinu pada $(0,1)$, tetapi konvergensinya tidak
seragam.
\end{exercise}

\begin{exnote}
Catatan: Dalam latihan sebelumnya, $(0,1)$ dipilih demi kesederhanaan.  Untuk latihan yang
lebih menantang, gantilah $(0,1)$ dengan $[0,1]$.
\end{exnote}

\begin{exercise}
Benar/Salah; buktikan atau berikan contoh penyangkal untuk pernyataan berikut:
Jika $\{ f_n \}_{n=1}^\infty$ merupakan barisan fungsi yang diskontinu di setiap titik pada $[0,1]$
dan konvergen seragam menuju fungsi $f$, maka $f$ diskontinu di setiap
titik.
\end{exercise}

\begin{exercise} \label{c1uniflim:exercise}
Untuk fungsi $f \colon [a,b] \to \R$ yang diferensiabel secara kontinu, definisikan
\begin{equation*}
\snorm{f}_{C^1} \coloneqq \snorm{f}_{[a,b]} + \snorm{f'}_{[a,b]} .
\end{equation*}
Andaikan $\{ f_n \}_{n=1}^\infty$ merupakan barisan fungsi yang diferensiabel secara
kontinu sedemikian sehingga untuk setiap $\epsilon >0$, terdapat $M$
sedemikian sehingga untuk semua $n,k \geq M$ berlaku
\begin{equation*}
\snorm{f_n-f_k}_{C^1} < \epsilon .
\end{equation*}
Tunjukkan bahwa $\{ f_n \}_{n=1}^\infty$ konvergen seragam menuju suatu fungsi
$f \colon [a,b] \to \R$ yang diferensiabel secara kontinu.
\end{exercise}

\begin{exnote}
Andaikan
$f \colon [0,1] \to \R$ terintegralkan Riemann.
Untuk dua latihan berikutnya, definisikan
bilangan
\begin{equation*}
\snorm{f}_{L^1} \coloneqq 
\int_0^1 \babs{f(x)}\,dx .
\end{equation*}
Jika $f$ terintegralkan, maka $\sabs{f}$ juga terintegralkan; lihat
\exerciseref{exercise:hardabsint}.
Bilangan tersebut disebut \emph{norma $L^1$}\index{L1-norm@norma $L^1$} dan
mendefinisikan jenis konvergensi lain yang sangat umum,
yang disebut
\emph{konvergensi $L^1$}\index{L1-convergence@konvergensi $L^1$}.
Namun, konsep ini sedikit lebih rumit.
\end{exnote}

\begin{exercise}
Andaikan $\{ f_n \}_{n=1}^\infty$ merupakan barisan fungsi yang terintegralkan Riemann pada $[0,1]$
dan konvergen seragam
ke $0$.  Tunjukkan bahwa
\begin{equation*}
\lim_{n\to\infty} \snorm{f_n}_{L^1} = 0 .
\end{equation*}
\end{exercise}

\begin{exercise}
Tentukan barisan $\{ f_n \}_{n=1}^\infty$ dari fungsi-fungsi terintegralkan Riemann
pada $[0,1]$ yang konvergen titik demi titik ke~$0$, tetapi sedemikian sehingga
\begin{equation*}
\lim_{n\to\infty} \snorm{f_n}_{L^1} \text{ tidak ada (bernilai } \infty\text{)}.
\end{equation*}
\end{exercise}

\begin{exercise}[Sulit] \label{exercise:dinisthm}
Buktikan \emph{\myindex{teorema Dini}}:
Misalkan $f_n \colon [a,b] \to \R$ merupakan barisan fungsi kontinu sedemikian sehingga
\begin{equation*}
0 \leq f_{n+1}(x) \leq f_n(x) \leq \cdots \leq f_1(x) 
\qquad \text{untuk semua } n \in \N \text{ dan } x \in [a,b].
\end{equation*}
Andaikan $\{ f_n \}_{n=1}^\infty$ konvergen titik demi titik ke $0$.
Tunjukkan bahwa $\{ f_n \}_{n=1}^\infty$ konvergen seragam ke nol.
\end{exercise}

\begin{exercise}
Andaikan $f_n \colon [a,b] \to \R$ merupakan barisan fungsi
kontinu yang
konvergen titik demi titik
menuju fungsi kontinu $f \colon [a,b] \to \R$.  Andaikan bahwa
untuk setiap $x \in [a,b]$,
barisan $\bigl\{ \babs{f_n(x)-f(x)} \bigr\}_{n=1}^\infty$ monoton.
Tunjukkan bahwa barisan $\{f_n\}_{n=1}^\infty$ konvergen seragam.
\end{exercise}

\begin{exercise}
\pagebreak[3]
Tentukan barisan fungsi terintegralkan Riemann $f_n \colon [0,1] \to \R$ sedemikian
sehingga $\{ f_n \}_{n=1}^\infty$ konvergen titik demi titik ke nol, dan sedemikian sehingga
\begin{enumerate}[a)]
\item
$\bigl\{ \int_0^1 f_n \bigr\}_{n=1}^\infty$ meningkat tanpa batas,
\item
$\bigl\{ \int_0^1 f_n \bigr\}_{n=1}^\infty$ merupakan barisan $-1,1,-1,1,-1,1, \ldots$.
\end{enumerate}
\end{exercise}

\begin{exnote}
Kita dapat mendefinisikan
\emph{\myindex{limit bersama}} dari barisan ganda
$\{ x_{n,m} \}_{n,m=1}^\infty$ yang terdiri atas bilangan
real (yakni suatu fungsi dari $\N \times \N$ ke $\R$).
Kita katakan bahwa $L$ merupakan limit bersama dari $\{ x_{n,m} \}_{n,m=1}^\infty$ dan menulis
\begin{equation*}
\lim_{\substack{n\to\infty\\m\to\infty}}
x_{n,m} = L ,
\qquad
\text{atau}
\qquad
\lim_{(n,m) \to \infty}
x_{n,m} = L ,
\end{equation*}
apabila untuk setiap $\epsilon > 0$, terdapat
$M$ sedemikian sehingga jika $n \geq M$ dan $m \geq M$, maka
$\sabs{x_{n,m} - L} < \epsilon$.
\end{exnote}

\begin{exercise}
Andaikan limit bersama (lihat di atas) dari $\{ x_{n,m} \}_{n,m=1}^\infty$ adalah $L$, dan andaikan
bahwa untuk semua $n$, $\lim\limits_{m \to \infty} x_{n,m}$ ada,
serta untuk semua $m$, $\lim\limits_{n \to \infty} x_{n,m}$ ada.  Tunjukkan bahwa
$\lim\limits_{n\to\infty}\lim\limits_{m \to \infty} x_{n,m}
=
\lim\limits_{m\to\infty}\lim\limits_{n \to \infty} x_{n,m} = L$.
\end{exercise}

\begin{exercise}
Adanya limit bersama (lihat di atas) tidak berarti bahwa limit berulang ada.
Pertimbangkan $x_{n,m} \coloneqq \frac{{(-1)}^{n+m}}{\min \{n,m \}}$.
\begin{enumerate}[a)]
\item
Tunjukkan bahwa untuk setiap $n$,
$\lim\limits_{m \to \infty} x_{n,m}$ tidak ada, dan untuk setiap $m$,
$\lim\limits_{n \to \infty} x_{n,m}$ juga tidak ada.  Jadi, baik
$\lim\limits_{n\to\infty}\lim\limits_{m \to \infty} x_{n,m}$ maupun
$\lim\limits_{m\to\infty}\lim\limits_{n \to \infty} x_{n,m}$ sama sekali tidak
bermakna.
\item
Tunjukkan bahwa limit bersama dari $\{ x_{n,m} \}_{n,m=1}^\infty$ ada dan sama dengan 0.
\end{enumerate}
\end{exercise}

\begin{exercise}
Kita katakan bahwa barisan fungsi $f_n \colon \R \to \R$
\emph{\myindex{konvergen seragam pada himpunan bagian kompak}}
\index{konvergensi seragam pada himpunan bagian kompak}
apabila untuk setiap $k \in \N$,
barisan $\{ f_n \}_{n=1}^\infty$ konvergen seragam pada $[-k,k]$.
\begin{enumerate}[a)]
\item
Buktikan bahwa jika $f_n \colon \R \to \R$ merupakan barisan
fungsi kontinu yang konvergen seragam pada himpunan bagian kompak, maka
limitnya kontinu.
\item 
Buktikan bahwa jika $f_n \colon \R \to \R$ merupakan barisan
fungsi yang terintegralkan Riemann pada setiap interval tertutup dan terbatas $[a,b]$,
dan konvergen seragam pada himpunan bagian kompak menuju $f \colon \R \to \R$,
maka untuk setiap interval $[a,b]$ berlaku $f \in \sR\bigl([a,b]\bigr)$, dan
$\int_a^b f = \lim\limits_{n\to\infty} \int_a^b f_n$.
\end{enumerate}
\end{exercise}

\begin{exercise}[Menantang]
Tentukan barisan fungsi kontinu $f_n \colon [0,1] \to \R$ yang
konvergen menuju fungsi popcorn $f \colon [0,1] \to \R$, yakni
fungsi sedemikian sehingga $f(\nicefrac{p}{q}) \coloneqq \nicefrac{1}{q}$ (jika $\nicefrac{p}{q}$
dalam bentuk paling sederhana) dan $f(x) \coloneqq 0$ jika $x$ irasional (perhatikan
bahwa $f(0) = f(1) = 1$);
lihat \exampleref{popcornfunction:example}.
Jadi, limit titik demi titik dari fungsi-fungsi kontinu dapat memiliki himpunan
diskontinuitas yang rapat.  Lihat juga latihan berikutnya.
\end{exercise}

\begin{exercise}[Menantang]
\myindex{Fungsi Dirichlet}
$f \colon [0,1] \to \R$, yakni
fungsi sedemikian sehingga $f(x) \coloneqq 1$ jika $x \in \Q$
dan $f(x) \coloneqq 0$ jika $x \notin \Q$,
bukan merupakan limit titik demi titik dari
fungsi-fungsi kontinu, walaupun hal ini sulit dibuktikan.
Namun, buktikan bahwa $f$ merupakan limit titik demi titik dari fungsi-fungsi yang masing-masing
merupakan limit titik demi titik dari
fungsi-fungsi kontinu.
\end{exercise}

\begin{exercise}
\leavevmode
\begin{enumerate}[a)]
\item
Tentukan barisan fungsi kontinu Lipschitz pada $[0,1]$
yang limit seragamnya adalah $\sqrt{x}$, yang bukan fungsi Lipschitz.
\item
Sebaliknya, tunjukkan bahwa jika $f_n \colon S \to \R$ merupakan fungsi Lipschitz
dengan konstanta seragam $K$ (artinya semuanya memenuhi definisi
dengan konstanta yang sama) dan $\{ f_n \}_{n=1}^\infty$ konvergen titik demi titik
menuju $f \colon S \to \R$,
maka limitnya, $f$, merupakan fungsi kontinu Lipschitz
dengan konstanta Lipschitz $K$.
\end{enumerate}
\end{exercise}

\begin{exercise}[memerlukan \sectionref{sec:moreonseries}]
Jika $\sum_{n=0}^\infty c_n {(x-a)}^n$ memiliki jari-jari konvergensi $\rho$,
tunjukkan bahwa integral suku demi sukunya,
$\sum_{n=1}^\infty \frac{c_{n-1}}{n} {(x-a)}^n$, memiliki jari-jari konvergensi
$\rho$.  Perhatikan bahwa di atas kita hanya membuktikan bahwa jari-jari konvergensinya
sekurang-kurangnya~$\rho$.
\end{exercise}

\begin{exercise}[memerlukan \sectionref{sec:moreonseries} dan \sectionref{sec:taylor}]
\pagebreak[3]
Andaikan $f(x) \coloneqq \sum_{n=0}^\infty c_n {(x-a)}^n$ konvergen pada $(a-\rho,a+\rho)$.
\begin{enumerate}[a)]
\item
Andaikan bahwa $f^{(k)}(a) = 0$ untuk semua $k=0,1,2,3,\ldots$.  Buktikan bahwa
$c_n = 0$ untuk semua $n$, atau dengan kata lain,
$f(x) = 0$ untuk semua $x \in (a-\rho,a+\rho)$.
\item
Dengan menggunakan bagian a), buktikan suatu versi dari apa yang disebut
\myquote{teorema identitas untuk fungsi analitik}: Jika terdapat $\epsilon > 0$
sedemikian sehingga $f(x) = 0$ untuk semua $x \in (a-\epsilon, a+\epsilon)$, maka
$f(x) = 0$ untuk semua $x \in (a-\rho,a+\rho)$.
\end{enumerate}
\end{exercise}

\begin{exercise}
Misalkan $f_n(x) \coloneqq \frac{x}{1+{(nx)}^2}$.  Perhatikan bahwa $f_n$ merupakan
fungsi-fungsi diferensiabel.
\begin{enumerate}[a)]
\item
Tunjukkan bahwa $\{ f_n \}_{n=1}^\infty$ konvergen seragam ke 0.
\item
Tunjukkan bahwa $\sabs{f_n'(x)} \leq 1$ untuk semua $x$ dan semua $n$.
\item
Tunjukkan bahwa $\{ f_n' \}_{n=1}^\infty$ konvergen titik demi titik menuju fungsi yang diskontinu di
titik asal.
\item
Misalkan $\{ a_n \}_{n=1}^\infty$ merupakan enumerasi bilangan rasional.
Definisikan
\begin{equation*}
g_n(x) \coloneqq \sum_{k=1}^n 2^{-k} f_n(x-a_k) .
\end{equation*}
Tunjukkan bahwa $\{ g_n \}_{n=1}^\infty$ konvergen seragam ke 0.
\item
Tunjukkan bahwa $\{ g_n' \}_{n=1}^\infty$ konvergen titik demi titik menuju fungsi $\psi$ yang
diskontinu di setiap bilangan rasional dan kontinu di setiap
bilangan irasional.  Secara khusus, $\lim\limits_{n\to\infty} g_n'(x) \neq 0$ untuk
setiap bilangan rasional~$x$.
\end{enumerate}
\end{exercise}

\begin{exercise}[memerlukan \sectionref{sec:impropriemann}]
Tunjukkan bahwa konvergensi seragam tidak cukup untuk mempertukarkan limit dengan
integral tak wajar pada interval tak terbatas.
Yakni, tentukan barisan fungsi $f_n \colon \R \to
\R$ yang terintegralkan Riemann pada setiap interval terbatas,
konvergen seragam ke nol, dan sedemikian sehingga
$\int_{-\infty}^\infty f_n = 1$ untuk setiap $n$.
\end{exercise}

%%%%%%%%%%%%%%%%%%%%%%%%%%%%%%%%%%%%%%%%%%%%%%%%%%%%%%%%%%%%%%%%%%%%%%%%%%%%%%

\sectionnewpage
\section{Teorema Picard}
\label{sec:picard}

%mbxINTROSUBSECTION

\sectionnotes{1--2 pertemuan kuliah (dapat dilewati tanpa masalah)}

Mata kuliah analisis semester pertama semestinya memiliki
teorema sekelas \emph{pi\`ece de r\'esistance}.
Kita memilih teorema yang pembuktiannya memadukan semua yang telah kita
pelajari.  Teorema ini lebih canggih daripada teorema dasar kalkulus,
teorema unggulan pertama dalam mata kuliah ini.  Teorema
yang kita maksud adalah teorema
Picard%
\footnote{Dinamai menurut matematikawan Prancis
\href{https://en.wikipedia.org/wiki/\%C3\%89mile_Picard}{Charles \'Emile Picard}
(1856--1941).}
tentang keberadaan dan ketunggalan solusi suatu persamaan diferensial biasa.
Pernyataan maupun pembuktiannya merupakan contoh indah dari apa yang dapat dilakukan
dengan materi yang sejauh ini telah kita kuasai.  Teorema ini juga merupakan contoh baik tentang penerapan analisis,
karena persamaan diferensial sangat diperlukan dalam berbagai cabang
sains.

\subsection{Persamaan diferensial biasa orde pertama}

Sains modern dijelaskan dalam bahasa
\emph{persamaan diferensial}\index{persamaan diferensial}.
Yakni, persamaan yang tidak hanya melibatkan fungsi yang belum diketahui, tetapi juga
turunannya.  Bentuk taktrivial paling sederhana dari persamaan diferensial adalah
\emph{\myindex{persamaan diferensial biasa orde pertama}}
\index{persamaan diferensial biasa}
\begin{equation*}
y' = F(x,y) .
\end{equation*}
Secara umum, kita juga menetapkan \emph{\myindex{syarat awal}}
$y(x_0)=y_0$.  Solusi
persamaan tersebut adalah fungsi $y(x)$ sedemikian sehingga
$y(x_0)=y_0$ dan $y'(x) = F\bigl(x,y(x)\bigr)$.
\figureref{fig:sampleslopes} menampilkan
\emph{\myindex{medan kemiringan}}, yaitu representasi grafis persamaan ini.
\begin{myfigureht}
\myincludegraphics{sampleslopes}{%
Sebagian bidang xy di sekitar titik asal dengan
kisi persegi panjang yang menunjukkan arah kemiringan.  Pada separuh kiri gambar,
kemiringannya mengarah ke bawah, tetapi semakin mendatar ketika bergerak ke atas.
Pada separuh kanan gambar, kemiringannya mengarah ke atas dan
kembali semakin mendatar ketika bergerak ke atas.
Sebuah titik (x subskrip nol, y subskrip nol) diberikan di kuadran pertama.
Grafik suatu solusi digambar mulai dari atas
sumbu x, turun, kemudian naik kembali dan melewati
titik yang diberikan.}
\caption{Sebuah \emph{medan kemiringan} yang memberikan kemiringan $F(x,y)$ di setiap titik;
dalam kasus ini $F(x,y)=x(1-y)$.  Sebuah solusi digambar melewati titik
$(x_0,y_0)=(1,0.3)$; perhatikan bagaimana grafik itu mengikuti arah-arah kemiringan tersebut.\label{fig:sampleslopes}}
\end{myfigureht}

Ketika $F$ hanya melibatkan variabel $x$, solusinya diberikan oleh
teorema dasar kalkulus.  Sebaliknya, ketika $F$ bergantung
pada $x$ maupun $y$, kita memerlukan perangkat yang jauh lebih kuat.  Solusi
tidak selalu ada, dan kalaupun ada, belum tentu merupakan satu-satunya solusi.
Teorema Picard memberikan syarat cukup tertentu bagi keberadaan
dan ketunggalan.

\subsection{Teorema tersebut}

Kita memerlukan definisi kekontinuan dalam dua variabel.  Suatu titik pada
bidang $\R^2 = \R \times \R$ dinyatakan dengan pasangan terurut $(x,y)$.
Demi kesederhanaan,
kita berikan definisi sekuensial kekontinuan berikut.

\begin{defn}
Misalkan $U \subset \R^2$ merupakan himpunan, $F \colon U \to \R$ suatu fungsi,
dan $(x,y) \in U$ suatu titik.  Fungsi $F$ disebut
\emph{kontinu di $(x,y)$}\index{fungsi kontinu dua variabel}
apabila untuk setiap barisan
$\bigl\{ (x_n,y_n) \bigr\}_{n=1}^\infty$ dari titik-titik di $U$ sedemikian sehingga
$\lim_{n\to\infty} x_n = x$ dan
$\lim_{n\to\infty} y_n = y$, berlaku
\begin{equation*}
\lim_{n \to \infty} F(x_n,y_n) = 
F(x,y) .
\end{equation*}
Kita katakan bahwa $F$ \emph{kontinu} apabila kontinu di semua titik dalam $U$.
\end{defn}

\begin{thm}[Teorema Picard tentang keberadaan dan ketunggalan]%
\index{teorema keberadaan dan ketunggalan}\index{teorema Picard}%
\label{thm:fs:picard}
Misalkan $I, J \subset \R$ merupakan interval tertutup dan terbatas,
misalkan $I^\circ$ dan $J^\circ$ merupakan interiornya%
\footnote{Yang dimaksud dengan interior $[a,b]$ adalah $(a,b)$.},
dan
misalkan $(x_0,y_0) \in I^\circ \times J^\circ$.
Andaikan $F \colon I \times J \to \R$ kontinu
dan Lipschitz dalam variabel kedua, yakni terdapat
$L \in \R$ sedemikian sehingga
\begin{equation*}
\babs{F(x,y) - F(x,z)} \leq L \sabs{y-z}
\qquad \text{untuk semua } y,z \in J, x \in I .
\end{equation*}
Maka terdapat $h > 0$ sedemikian sehingga $[x_0-h,x_0+h] \subset I$ dan terdapat tepat satu fungsi
diferensiabel $f \colon [x_0 - h, x_0 + h] \to J \subset \R$ sedemikian sehingga
\begin{equation} \label{picard:diffeq}
f'(x) = F\bigl(x,f(x)\bigr) \qquad \text{dan} \qquad f(x_0) = y_0.
\end{equation}
\end{thm}

\begin{proof}
\pagebreak[2]
Andaikan kita dapat menemukan solusi $f$.  Dengan menggunakan teorema dasar
kalkulus, kita mengintegralkan persamaan
$f'(x) = F\bigl(x,f(x)\bigr)$, $f(x_0) = y_0$, dan menuliskan \eqref{picard:diffeq}
sebagai persamaan integral
\begin{equation} \label{picard:inteq}
f(x) = y_0 + \int_{x_0}^x F\bigl(t,f(t)\bigr)\,dt .
\end{equation}
Gagasan pembuktian kita adalah memasukkan hampiran-hampiran suatu
solusi ke ruas kanan \eqref{picard:inteq} untuk memperoleh hampiran yang lebih baik pada
ruas kiri \eqref{picard:inteq}.  Kita berharap bahwa pada akhirnya barisan tersebut
konvergen dan limitnya memenuhi
\eqref{picard:inteq}, dan karenanya juga \eqref{picard:diffeq}.
Teknik berikut disebut \emph{\myindex{iterasi Picard}},
dan masing-masing fungsi $f_k$ disebut
\emph{iterat Picard}\index{iterat Picard}.

Tanpa mengurangi keumuman, andaikan $x_0 = 0$ (lihat latihan di bawah).
Latihan lainnya
menyatakan bahwa $F$ terbatas karena kontinu.
Oleh karena itu, pilih $M > 0$ sedemikian sehingga
$\babs{F(x,y)} \leq M$ untuk semua $(x,y) \in I\times J$.
Pilih $\alpha > 0$ sedemikian sehingga
$[-\alpha,\alpha] \subset I$ dan $[y_0-\alpha, y_0 + \alpha] \subset J$.
Definisikan
\begin{equation*}
h \coloneqq \min \left\{ \alpha, \frac{\alpha}{M+L\alpha} \right\} .
\end{equation*}
Perhatikan bahwa
$[-h,h] \subset I$.

Tetapkan $f_0(x) \coloneqq y_0$.
Kita definisikan $f_k$ secara induktif.
Dengan mengandaikan $f_{k-1}([-h,h]) \subset [y_0-\alpha,y_0+\alpha]$,
kita melihat bahwa
$F\bigl(t,f_{k-1}(t)\bigr)$
terdefinisi dengan baik sebagai fungsi terhadap $t$ untuk $t \in [-h,h]$.
Jika $f_{k-1}$ kontinu
pada $[-h,h]$, maka
$F\bigl(t,f_{k-1}(t)\bigr)$
kontinu sebagai
fungsi terhadap $t$ pada $[-h,h]$ (diserahkan sebagai latihan).
Definisikan
\begin{equation*}
f_k(x) \coloneqq y_0+ \int_{0}^x F\bigl(t,f_{k-1}(t)\bigr)\,dt ,
\end{equation*}
dan $f_k$ kontinu pada $[-h,h]$ berdasarkan teorema dasar kalkulus.
Untuk melihat bahwa $f_k$ memetakan $[-h,h]$ ke $[y_0-\alpha,y_0+\alpha]$, kita hitung, untuk
$x \in [-h,h]$,
\begin{equation*}
\sabs{f_k(x) - y_0} = 
\abs{\int_{0}^x F\bigl(t,f_{k-1}(t)\bigr)\,dt }
\leq
M\sabs{x}
\leq
Mh
\leq
M
\frac{\alpha}{M+L\alpha}
\leq \alpha .
\end{equation*}
Selanjutnya kita definisikan $f_{k+1}$ menggunakan $f_k$, dan demikian seterusnya.
Dengan demikian, secara induktif kita telah mendefinisikan barisan $\{ f_k \}_{k=1}^\infty$
dari fungsi-fungsi kontinu.
Kita perlu menunjukkan bahwa barisan tersebut konvergen menuju fungsi $f$ yang memenuhi
persamaan~\eqref{picard:inteq}, dan karenanya~\eqref{picard:diffeq}.

Kita ingin menunjukkan bahwa barisan $\{ f_k \}_{k=1}^\infty$ konvergen seragam
menuju suatu fungsi pada $[-h,h]$.  Pertama-tama, untuk $t \in [-h,h]$,
kita memiliki batas berguna
berikut
\begin{equation*}
\Babs{F\bigl(t,f_{n}(t)\bigr) - 
F\bigl(t,f_{k}(t)\bigr)}
\leq
L \babs{f_n(t)-f_k(t)}
\leq
L \snorm{f_n-f_k}_{[-h,h]} ,
\end{equation*}
di mana $\snorm{f_n-f_k}_{[-h,h]}$ adalah norma seragam,
yakni supremum dari $\babs{f_n(t)-f_k(t)}$ untuk $t \in [-h,h]$.
Sekarang perhatikan bahwa $\sabs{x} \leq h \leq \frac{\alpha}{M+L\alpha}$.
Oleh karena itu,
\begin{equation*}
\begin{split}
\babs{f_n(x) - f_k(x)}
& =
\abs{\int_{0}^x F\bigl(t,f_{n-1}(t)\bigr)\,dt 
-
\int_{0}^x F\bigl(t,f_{k-1}(t)\bigr)\,dt}
\\
& =
\abs{\int_{0}^x
\Bigl(
F\bigl(t,f_{n-1}(t)\bigr)-
F\bigl(t,f_{k-1}(t)\bigr)
\Bigr)
\,dt}
\\
& \leq
L
\,
\snorm{f_{n-1}-f_{k-1}}_{[-h,h]}
\,
\sabs{x}
\\
& \leq
\frac{L\alpha}{M+L\alpha}
\snorm{f_{n-1}-f_{k-1}}_{[-h,h]} .
\end{split}
\end{equation*}
Misalkan $C \coloneqq \frac{L\alpha}{M+L\alpha}$ dan perhatikan bahwa $C < 1$.
Dengan mengambil supremum pada ruas kiri, kita peroleh
\begin{equation*}
\snorm{f_n-f_k}_{[-h,h]} \leq C \snorm{f_{n-1}-f_{k-1}}_{[-h,h]} .
\end{equation*}
Tanpa mengurangi keumuman,
andaikan $n \geq k$.  Dengan \hyperref[induction:thm]{induksi}, kita dapat menunjukkan bahwa
\begin{equation*}
\snorm{f_n-f_k}_{[-h,h]} \leq C^{k} \snorm{f_{n-k}-f_{0}}_{[-h,h]} .
\end{equation*}
Untuk $x \in [-h,h]$, berlaku
\begin{equation*}
\babs{f_{n-k}(x)-f_{0}(x)}
=
\babs{f_{n-k}(x)-y_0}
\leq \alpha .
\end{equation*}
Oleh karena itu,
\begin{equation*}
\snorm{f_n-f_k}_{[-h,h]} \leq C^{k} \snorm{f_{n-k}-f_{0}}_{[-h,h]} \leq C^{k} \alpha .
\end{equation*}
Karena $C < 1$, $\{f_n\}_{n=1}^\infty$ bersifat Cauchy seragam.
Berdasarkan \propref{prop:uniformcauchy}, $\{ f_n \}_{n=1}^\infty$
konvergen seragam pada $[-h,h]$ menuju suatu fungsi $f \colon [-h,h] \to \R$.
Fungsi $f$ merupakan limit seragam dari fungsi-fungsi kontinu, dan karenanya
kontinu.  Selain itu, karena $f_n\bigl([-h,h]\bigr) \subset
[y_0-\alpha,y_0+\alpha]$ untuk semua $n$,
maka $f\bigl([-h,h]\bigr) \subset [y_0-\alpha,y_0+\alpha]$
(mengapa?).


Sekarang kita perlu menunjukkan bahwa $f$ memenuhi \eqref{picard:inteq}.
Pertama-tama, seperti sebelumnya, kita perhatikan
\begin{equation*}
\Babs{F\bigl(t,f_{n}(t)\bigr) - 
F\bigl(t,f(t)\bigr)}
\leq
L \babs{f_n(t)-f(t)}
\leq
L \snorm{f_n-f}_{[-h,h]} .
\end{equation*}
Karena $\snorm{f_n-f}_{[-h,h]}$ konvergen ke nol,
$F\bigl(t,f_n(t)\bigr)$ konvergen seragam ke $F\bigl(t,f(t)\bigr)$
untuk $t \in [-h,h]$.  Oleh karena itu, untuk $x \in [-h,h]$,
konvergensinya seragam %(mengapa?)
untuk $t \in [0,x]$ (atau $[x,0]$ jika $x < 0$).  Dengan demikian,
\begin{align*}
y_0
+
\int_0^{x}
F\bigl(t,f(t)\bigr)\,dt
& =
y_0
+
\int_0^{x}
F\bigl(t,\lim_{n\to\infty} f_n(t)\bigr)\,dt
& &
\\
& =
y_0
+
\int_0^{x}
\lim_{n\to\infty} F\bigl(t,f_n(t)\bigr)\,dt
& & \text{(berdasarkan kekontinuan } F\text{)}
\\
& =
\lim_{n\to\infty} 
\left(
y_0
+
\int_0^{x}
F\bigl(t,f_n(t)\bigr)\,dt
\right)
& & \text{(berdasarkan konvergensi seragam)}
\\
& =
\lim_{n\to\infty} 
f_{n+1}(x)
=
f(x) .
& &
\end{align*}
Kita menerapkan teorema dasar kalkulus (\thmref{thm:FTCv2}) untuk menunjukkan bahwa
$f$ diferensiabel dan turunannya adalah $F\bigl(x,f(x)\bigr)$.  Jelas bahwa
$f(0) = y_0$.

Terakhir, kita tinggal menunjukkan ketunggalan.  Andaikan $g \colon [-h,h]
\to J \subset \R$ merupakan solusi lain.
Seperti sebelumnya, kita gunakan fakta bahwa
$\babs{F\bigl(t,f(t)\bigr) - F\bigl(t,g(t)\bigr)} \leq L \snorm{f-g}_{[-h,h]}$.
Maka
\begin{equation*}
\begin{split}
\babs{f(x)-g(x)}
& =
\abs{
y_0
+
\int_0^{x}
F\bigl(t,f(t)\bigr)\,dt
-
\left(
y_0
+
\int_0^{x}
F\bigl(t,g(t)\bigr)\,dt
\right)
}
\\
& =
\abs{
\int_0^{x}
\Bigl(
F\bigl(t,f(t)\bigr)
-
F\bigl(t,g(t)\bigr)
\Bigr)
\,dt
}
\\
& \leq
L\snorm{f-g}_{[-h,h]}\sabs{x}
\leq
Lh\snorm{f-g}_{[-h,h]}
\leq
\frac{L\alpha}{M+L\alpha}\snorm{f-g}_{[-h,h]} .
\end{split}
\end{equation*}
Seperti
sebelumnya, $C = \frac{L\alpha}{M+L\alpha} < 1$.  Dengan mengambil supremum terhadap $x \in
[-h,h]$ pada
ruas kiri, kita peroleh
\begin{equation*}
\snorm{f-g}_{[-h,h]} \leq C \snorm{f-g}_{[-h,h]} .
\end{equation*}
Hal ini hanya mungkin jika $\snorm{f-g}_{[-h,h]} = 0$.  Oleh karena itu, $f=g$, dan
solusinya tunggal.
\end{proof}

\subsection{Contoh}

Mari kita lihat beberapa contoh.  Pembuktian teorema tersebut
memberikan cara eksplisit untuk menemukan $h$ yang sesuai.  Namun, pembuktian itu tidak
memberikan $h$ terbaik.  Sering kali kita dapat menemukan $h$ yang jauh lebih besar
dan tetap memenuhi kesimpulan teorema.

Pembuktian tersebut juga memberikan iterat-iterat Picard sebagai hampiran terhadap
solusi.  Jadi, pembuktian itu sesungguhnya memberi tahu kita cara memperoleh
solusi, bukan sekadar bahwa solusinya ada.

\begin{example} \label{example:picardexponential}
Pertimbangkan
\begin{equation*}
f'(x) = f(x), \qquad f(0) = 1 .
\end{equation*}
Yakni, kita andaikan $F(x,y) = y$, dan kita mencari fungsi
$f$ sedemikian sehingga $f'(x) = f(x)$.
Untuk sementara, lupakan bahwa kita telah menyelesaikan persamaan ini dalam
\sectionref{sec:logandexp}.  Lihat juga \figureref{fig:exp} untuk plot yang sekaligus memperlihatkan
persamaan tersebut, dengan
kemiringan $F(x,y)=y$ di setiap titik, dan solusinya,
yaitu fungsi eksponensial yang memenuhi $f(0)=1$.

Kita pilih sebarang $I$ yang memuat 0
di interiornya.
Kita pilih sebarang $J$ yang memuat 1 di interiornya.  Kita dapat
menggunakan $L = 1$.
Teorema tersebut menjamin adanya $h > 0$ sedemikian sehingga
terdapat solusi tunggal $f \colon [-h,h] \to \R$.  Solusi ini
biasanya dinotasikan dengan
\begin{equation*}
e^x \coloneqq f(x) .
\end{equation*}
Kami serahkan kepada pembaca untuk memeriksa bahwa dengan memilih $I$ dan $J$
cukup besar, pembuktian teorema menjamin bahwa
kita dapat memilih $\alpha$ sehingga memperoleh sebarang
$h$ yang kita inginkan asalkan $h < \nicefrac{1}{2}$.  Perhitungannya kita hilangkan.
Tentu saja, kita tahu %(though we have not proved)
bahwa fungsi ini ada
untuk semua $x$, sehingga sebarang $h$ semestinya dapat digunakan, tetapi
pembuktian teorema hanya memberikan $h < \nicefrac{1}{2}$.

Dengan penalaran yang sama seperti di atas,
apa pun nilai $x_0$ dan $y_0$,
pembuktian tersebut menjamin sebarang $h$ asalkan $h < \nicefrac{1}{2}$.
Tetapkan $h$ semacam itu.
Kita memperoleh fungsi tunggal yang terdefinisi pada $[x_0-h,x_0+h]$.  Setelah mendefinisikan
fungsi pada $[-h,h]$, kita menemukan solusi pada interval $[0,2h]$
dan memperhatikan bahwa kedua fungsi harus berimpit pada $[0,h]$ karena ketunggalan.
Dengan demikian, secara iteratif kita membangun fungsi eksponensial untuk semua $x \in \R$.
Oleh karena itu, teorema Picard dapat digunakan untuk membuktikan keberadaan dan ketunggalan
fungsi eksponensial.

Mari kita hitung iterat-iterat Picard.
Kita mulai dengan fungsi konstan $f_0(x) \coloneqq 1$.  Maka
\begin{align*}
f_1(x) & = 1 + \int_0^x f_0(s)\,ds =
1+x, \\
f_2(x) & = 1 + \int_0^x f_1(s)\,ds =
1 + \int_0^x (1+s)\,ds = 1 + x + \frac{x^2}{2}, \\
f_3(x) & = 1 + \int_0^x f_2(s)\,ds =
1 + \int_0^x \left(1+ s + \frac{s^2}{2} \right)\,ds =
1 + x + \frac{x^2}{2} + \frac{x^3}{6} .
\end{align*}
Kita mengenali bagian awal deret Taylor untuk fungsi eksponensial.
Lihat \figureref{fig:exppicard}.
\begin{myfigureht}
\myincludegraphics{exppicardfig}{%
Grafik fungsi eksponensial berupa
kurva penuh pada
daerah dengan x di antara sedikit kurang dari minus 1 dan
sedikit lebih dari 1.
Beberapa kurva putus-putus memberikan hampiran.
Semuanya melewati titik (0,1).
Pertama, sebuah garis horizontal.  Lalu sebuah garis
dengan kemiringan yang sama seperti fungsi eksponensial di (0,1).
Kemudian dua hampiran melengkung lainnya yang semakin mendekati
fungsi eksponensial, terutama di dekat x sama dengan 0.}
\caption{Fungsi eksponensial (garis penuh) bersama dengan $f_0$, $f_1$, $f_2$,
$f_3$ (putus-putus).\label{fig:exppicard}}
\end{myfigureht}
\end{example}

\begin{example}
Pertimbangkan persamaan
\begin{equation*}
f'(x) = {\bigl(f(x)\bigr)}^2 \qquad \text{dan} \qquad f(0)=1.
\end{equation*}
Melalui persamaan diferensial elementer kita mengetahui bahwa
\begin{equation*}
f(x) = \frac{1}{1-x}
\end{equation*}
merupakan solusinya.
Solusi tersebut hanya terdefinisi pada $(-\infty,1)$.  Artinya,
kita dapat menggunakan $h < 1$, tetapi tidak pernah $h$ yang lebih besar.
Fungsi yang memetakan $y$ ke $y^2$
bukan merupakan fungsi Lipschitz pada seluruh $\R$.
Ketika kita mendekati $x=1$ dari kiri, nilai solusinya semakin
besar.  Turunan solusi bertambah sebagai $y^2$, sehingga
$L$ yang diperlukan harus semakin besar ketika $y_0$ membesar.
Jika kita menerapkan
teorema tersebut dengan $x_0$ dekat 1 dan $y_0 = \frac{1}{1-x_0}$, kita mendapati
bahwa $h$ yang dijamin oleh pembuktian semakin kecil ketika $x_0$
mendekati~1.

Nilai $h$ dari pembuktian bukanlah nilai $h$ terbaik.
Dengan memilih $\alpha$ secara tepat, pembuktian teorema menjamin
$h=1-\nicefrac{\sqrt{3}}{2} \approx 0.134$ (perhitungannya kita
hilangkan) untuk $x_0=0$ dan $y_0=1$, walaupun
di atas kita melihat bahwa sebarang $h < 1$ semestinya dapat digunakan.
\end{example}

\begin{example}
Pertimbangkan persamaan
\begin{equation*}
f'(x) = 2 \sqrt{\babs{f(x)}}, \qquad f(0) = 0 .
\end{equation*}
Fungsi $F(x,y) = 2 \sqrt{\sabs{y}}$ kontinu,
tetapi tidak Lipschitz dalam $y$ (mengapa?).
Persamaan tersebut tidak memenuhi hipotesis teorema.
Fungsi
\begin{equation*}
f(x) =
\begin{cases}
x^2 & \text{jika } x \geq 0,\\
-x^2 & \text{jika } x < 0,
\end{cases}
\end{equation*}
merupakan solusi, tetapi $f(x) = 0$ juga merupakan solusi.
Jadi, solusi ada, tetapi tidak tunggal.
\end{example}

\begin{example}
Pertimbangkan $y' = \varphi(x)$ dengan $\varphi(x) \coloneqq 0$ jika $x \in \Q$ dan
$\varphi(x)\coloneqq 1$ jika $x
\notin \Q$.  Dengan kata lain, $F(x,y) = \varphi(x)$ diskontinu.
Persamaan tersebut tidak memiliki solusi, apa pun syarat
awalnya.
Seandainya ada, solusi itu akan memiliki
turunan $\varphi$, tetapi $\varphi$ tidak memiliki sifat nilai antara
di titik mana pun (mengapa?).  Berdasarkan
\hyperref[thm:darboux]{teorema Darboux}, tidak ada solusi.
\end{example}

Contoh-contoh tersebut menunjukkan bahwa tanpa syarat Lipschitz, solusi mungkin
ada tetapi tidak tunggal, dan tanpa kekontinuan $F$, solusi mungkin
sama sekali tidak ada.
Sesungguhnya terdapat suatu teorema, yaitu teorema keberadaan Peano, yang menyatakan bahwa jika $F$
kontinu, maka solusi ada (tetapi mungkin tidak tunggal).

\begin{remark}
Kita dapat memperlemah apa yang dimaksud dengan \myquote{solusi dari $y' = F(x,y)$}
dengan berfokus pada persamaan integral
$f(x) = y_0 + \int_{x_0}^x F\bigl(t,f(t)\bigr) \, dt$.  Sebagai contoh,
misalkan $H$ merupakan
fungsi Heaviside%
\footnote{Dinamai menurut
insinyur, matematikawan, dan fisikawan Inggris
\href{https://en.wikipedia.org/wiki/Oliver_Heaviside}{Oliver Heaviside}
(1850--1925).},
yakni $H(x) \coloneqq 0$ untuk $x < 0$ dan $H(x) \coloneqq 1$ untuk $x \geq 0$.
Maka
$y' = H(x)$, $y(0) = 0$,
merupakan persamaan yang umum.
\myquote{Solusinya}
adalah fungsi ramp $f(x) \coloneqq 0$ jika $x < 0$ dan $f(x) \coloneqq x$ jika $x \geq 0$,
karena fungsi ini memenuhi
$f(x) = \int_0^x H(t)\, dt$.  Namun, perhatikan bahwa $f'(0)$ tidak ada,
sehingga $f$ hanya merupakan apa yang disebut \emph{\myindex{solusi lemah}} dari
persamaan diferensial tersebut.
\end{remark}


\subsection{Latihan}

\begin{exercise}
Misalkan $I, J \subset \R$ merupakan interval.
Misalkan $F \colon I \times J \to \R$ merupakan fungsi kontinu
dua variabel,
dan andaikan $f \colon I \to J$ merupakan fungsi kontinu.
Tunjukkan bahwa $F\bigl(x,f(x)\bigr)$ merupakan fungsi kontinu pada $I$.
\end{exercise}

\begin{exercise}
Misalkan $I, J \subset \R$ merupakan interval tertutup dan terbatas.
Tunjukkan bahwa jika $F \colon I \times J \to \R$ kontinu,
maka $F$ terbatas.
\end{exercise}

\begin{exercise}
Kita membuktikan teorema Picard dengan mengandaikan bahwa $x_0 = 0$.
Buktikan pernyataan lengkap teorema Picard untuk sebarang $x_0$.
\end{exercise}

\begin{exercise}
Misalkan persamaan kita adalah $f'(x)=x f(x)$.  Mulailah dengan syarat awal
$f(0)=2$ dan tentukan iterat-iterat Picard $f_0,f_1,f_2,f_3,f_4$.
\end{exercise}

\begin{exercise}
Andaikan $F \colon I \times J \to \R$
merupakan fungsi yang kontinu dalam variabel pertama,
yakni, untuk setiap $y$ yang tetap, fungsi yang memetakan $x$ ke $F(x,y)$
kontinu.  Selain itu, andaikan $F$ Lipschitz dalam variabel kedua,
yakni terdapat bilangan $L$ sedemikian sehingga
\begin{equation*}
\babs{F(x,y) - F(x,z)} \leq L \sabs{y-z}
\qquad \text{untuk semua } y,z \in J, x \in I .
\end{equation*}
Tunjukkan bahwa $F$ kontinu sebagai fungsi dua variabel.  Oleh karena itu,
hipotesis dalam teorema dapat dibuat lebih lemah lagi.
\end{exercise}

\begin{exercise}
\pagebreak[2]
Jenis persamaan yang umum adalah
\emph{\myindex{persamaan diferensial linear orde pertama}}, yakni
persamaan berbentuk
\begin{equation*}
y' + p(x) y = q(x) , \qquad y(x_0) = y_0 .
\end{equation*}
Buktikan teorema Picard untuk persamaan linear.  Andaikan $I$ merupakan
interval, $x_0 \in I$, serta $p \colon I \to \R$ dan $q \colon I \to \R$
kontinu.
Tunjukkan bahwa terdapat tepat satu fungsi diferensiabel $f \colon I \to \R$
sedemikian sehingga $y = f(x)$
memenuhi persamaan dan syarat awal tersebut.
Petunjuk: Andaikan fungsi eksponensial ada dan gunakan rumus faktor
integrasi untuk membuktikan keberadaan $f$ (buktikan bahwa rumusnya bekerja, lalu buktikan ketunggalannya):
\begin{equation*}
f(x) \coloneqq e^{-\int_{x_0}^x p(s)\, ds} \left( \int_{x_0}^x e^{\int_{x_0}^t p(s)\, ds}
q(t) \,dt + y_0 \right).
\end{equation*}
\end{exercise}

\begin{exercise}
Pertimbangkan persamaan $f'(x) = f(x)$
dari \exampleref{example:picardexponential}.  Tunjukkan bahwa untuk sebarang $x_0$,
sebarang $y_0$, dan sebarang bilangan positif $h < \nicefrac{1}{2}$, kita dapat memilih $\alpha >
0$ yang cukup besar sehingga pembuktian teorema Picard menjamin adanya solusi
dengan syarat awal $f(x_0) = y_0$ pada interval
$[x_0-h,x_0+h]$.
\end{exercise}

\begin{exercise}
\pagebreak[3]
Pertimbangkan persamaan $y' = y^{1/3}x$.
\begin{enumerate}[a)]
\item
Tunjukkan bahwa untuk syarat awal $y(1)=1$, teorema Picard berlaku.
Tentukan $\alpha > 0$, $M$, $L$, dan $h$ yang dapat digunakan dalam pembuktian.
\item
Tunjukkan bahwa untuk syarat awal $y(1) = 0$, teorema Picard
tidak berlaku.
\item
Meskipun demikian, tentukan solusi untuk $y(1) = 0$.
\end{enumerate}
\end{exercise}

\begin{exercise}
Pertimbangkan persamaan $x y' = 2y$.
\begin{enumerate}[a)]
\item
Tunjukkan bahwa $y = Cx^2$ merupakan solusi untuk setiap konstanta $C$.
\item
Tunjukkan bahwa untuk setiap $x_0 \neq 0$ dan setiap $y_0$, teorema
Picard berlaku bagi syarat awal $y(x_0) = y_0$.
\item
Tunjukkan bahwa $y(0) = y_0$ dapat diselesaikan jika dan hanya jika $y_0 = 0$.
\end{enumerate}
\end{exercise}
